\part{Analytic, probabilistic, and Mellin extensions}

\section*{Editorial map of the independent extensions}
\addcontentsline{toc}{section}{Editorial map of the independent extensions}
The finite geometric-comb algebra is stated once in the preceding part. This part retains only the independent material not subsumed by that spine: confluent modal filters, regular variation, geometric-knot splines, geometric-uniform laws, Mellin symbols, superflat divergence, and the reciprocal Rvachev filter.

\section{The Gaussian origin filter as a dilation polynomial}
\label{g1:sec:operator}

The preceding part derives the origin weights from geometric Lagrange
interpolation.  The following compact operator form is the common engine for
the extensions in this part.  For $0<q<1$ put
\begin{align}
 \lambda_{n,k}^{(q)}
 &:={(-1)^{n-k}q^{\binom{n-k+1}{2}}
       \over (q;q)_k(q;q)_{n-k}},
 \qquad 0\le k\le n,
 \label{g1:eq:lambda}\\
 (\Gc_{n,q}f)(A)&:=\sum_{k=0}^n\lambda_{n,k}^{(q)}f(Aq^k).
 \label{g1:eq:G-definition}
\end{align}

\begin{theorem}[Dilation-polynomial factorization]
\label{g1:thm:operator-factorization}
Let $(T_qf)(x)=f(qx)$.  Then
\begin{align}
 W_{n,q}(z):=\sum_{k=0}^n\lambda_{n,k}^{(q)}z^k
 &=\prod_{j=1}^n{z-q^j\over1-q^j},
 \label{g1:eq:row-polynomial}\\
 \Gc_{n,q}=W_{n,q}(T_q)
 &={1\over(q;q)_n}\prod_{j=1}^n(T_q-q^jI),
 \label{g1:eq:operator-factorization}\\
 \Gc_{0,q}=I,
 \qquad
 \Gc_{n,q}&={T_q-q^nI\over1-q^n}\Gc_{n-1,q}.
 \label{g1:eq:operator-recursion}
\end{align}
Consequently the filter can be evaluated in place: starting with
$y_k=f(Aq^k)$, at stage $j=1,\ldots,n$ replace
\begin{equation}
 y_k\longleftarrow{y_{k+1}-q^jy_k\over1-q^j}
 \qquad(0\le k\le n-j).
 \label{g1:eq:triangular-filter}
\end{equation}
The final entry is $(\Gc_{n,q}f)(A)$.
\end{theorem}

\begin{proof}
The finite $q$-binomial theorem expands the product in
\eqref{g1:eq:row-polynomial} and gives exactly the coefficients
\eqref{g1:eq:lambda}.  Substituting the commuting operator $T_q$ for $z$
proves the factorization; separating its last factor proves the recursion.
Applied to the sample vector, that recursion is precisely the in-place update
\eqref{g1:eq:triangular-filter}.
\end{proof}

\begin{corollary}[Exact power multiplier and all integer moments]
\label{g1:cor:power-multiplier}
For $A>0$ and a chosen branch of $A^\alpha,q^\alpha$,
\begin{equation}
 (\Gc_{n,q}x^\alpha)(A)
 =A^\alpha W_{n,q}(q^\alpha)
 =A^\alpha q^{\alpha n}{(q^{1-\alpha};q)_n\over(q;q)_n}.
 \label{g1:eq:pure-power}
\end{equation}
In particular, for every $d\in\NaturalNumbers$,
\begin{equation}
 \sum_{k=0}^n\lambda_{n,k}^{(q)}(Aq^k)^d
 =\begin{cases}
  1,&d=0,\\
  0,&1\le d\le n,\\
  (-1)^nA^dq^{\binom{n+1}{2}}
  \GaussianBinomial{d-1}{n}{q},&d\ge n+1.
 \end{cases}
 \label{g1:eq:all-moments}
\end{equation}
\end{corollary}

\begin{proof}
The function $x^\alpha$ is an eigenfunction of $T_q$ with eigenvalue
$q^\alpha$, so \eqref{g1:eq:row-polynomial} gives the first equality.
Factoring $q^\alpha$ from every numerator factor gives the Pochhammer form.
For integral $d$, its factors give the vanishing range; when $d\ge n+1$,
reversing the factors and using the factorial definition of the Gaussian
coefficient gives \eqref{g1:eq:all-moments}.
\end{proof}

\begin{theorem}[Exact analytic tail and reciprocal-$q$ majorant]
\label{g1:thm:sharp-majorant}
Let $f(z)=\sum_{m\ge0}a_mz^m$.  Formally, and analytically whenever the
series is absolutely convergent,
\begin{equation}
 (\Gc_{n,q}f)(A)-f(0)
 =(-1)^nq^{\binom{n+1}{2}}
  \sum_{r\ge0}\GaussianBinomial{n+r}{n}{q}a_{n+1+r}A^{n+1+r}.
 \label{g1:eq:exact-analytic-tail}
\end{equation}
If $|a_m|\le C R^{-m}$ and $t=|A|/R<1$, then
\begin{equation}
 \left|(\Gc_{n,q}f)(A)-f(0)\right|
 \le Cq^{\binom{n+1}{2}}{t^{n+1}\over(t;q)_{n+1}}.
 \label{g1:eq:sharp-majorant}
\end{equation}
This majorant is coefficientwise sharp: equality in the majorizing sum is
attained when the relevant coefficients have the common phase and magnitudes
$CR^{-m}$.
\end{theorem}

\begin{proof}
Apply \eqref{g1:eq:all-moments} term by term and reindex $m=n+1+r$.
Taking absolute values leaves the series
\[
 Cq^{\binom{n+1}{2}}t^{n+1}
 \sum_{r\ge0}\GaussianBinomial{n+r}{n}{q}t^r.
\]
The reciprocal finite $q$-binomial theorem identifies the last sum with
$(t;q)_{n+1}^{-1}$.  The asserted phase choice makes every triangle
inequality an equality.
\end{proof}

\section{Confluent modal extrapolation}
\label{g1:sec:modal}

Ordinary polynomial cancellation assumes an integer-power Taylor expansion.
Endpoint asymptotics in differential equations, Mellin analysis, and inverse
special functions often contain fractional powers and logarithms.  The
dilation-operator viewpoint gives the correct generalization immediately.

\begin{definition}[Mode polynomial]
Let $\alpha_1,\ldots,\alpha_J\in\ComplexNumbers$ have pairwise distinct spectral
points $q^{\alpha_j}$, none equal to $1$, and let $m_1,\ldots,m_J\in\NaturalNumbers$, and put $N=\sum_jm_j$.  Define
\begin{equation}
 W_{\bm\alpha,\bm m}(z)
 :=\prod_{j=1}^J
 \left(\frac{z-q^{\alpha_j}}{1-q^{\alpha_j}}\right)^{m_j}
 =\sum_{k=0}^Nw_kz^k,
 \label{g1:eq:mode-polynomial}
\end{equation}
and
\[
 (\Rc_{\bm\alpha,\bm m}f)(x):=\sum_{k=0}^Nw_kf(q^kx).
\]
\end{definition}

\begin{theorem}[Confluent annihilation]
The filter \eqref{g1:eq:mode-polynomial} preserves constants and, for every
$j$ and $0\le s<m_j$, annihilates
\[
 x^{\alpha_j}(\log x)^s.
\]
It is the unique filter of degree at most $N$ with these properties.
\end{theorem}

\begin{proof}
For a general polynomial $W(z)=\sum_kw_kz^k$,
\begin{align*}
 \sum_kw_k(q^kx)^\alpha(\log(q^kx))^s
 &=x^\alpha\sum_kw_kq^{k\alpha}
   (\log x+k\log q)^s\\
 &=x^\alpha(\log x+\partial_\alpha)^sW(q^\alpha).
\end{align*}
A zero of multiplicity $m_j$ at $q^{\alpha_j}$ makes the right side vanish for
$s<m_j$.  The normalization $W(1)=1$ preserves constants.  Conversely, the
annihilation equations imply zeros with the stated multiplicities; degree and
normalization then force \eqref{g1:eq:mode-polynomial}.
\end{proof}

\begin{corollary}[Exact general multiplier]
For any $\alpha\in\ComplexNumbers$ and $s\in\NaturalNumbers$,
\begin{equation}
 \Rc_{\bm\alpha,\bm m}
 [x^\alpha(\log x)^s]
 =x^\alpha(\log x+\partial_\alpha)^s
  W_{\bm\alpha,\bm m}(q^\alpha).
 \label{g1:eq:general-mode-multiplier}
\end{equation}
\end{corollary}

\begin{proof}
Because $x^\alpha(\log x)^s=\partial_\alpha^s x^\alpha$ and the
filter is a finite linear combination of dilations, differentiation in
$\alpha$ commutes with the filter.  Apply the pure-power multiplier and
expand $\partial_\alpha^s\!\left(x^\alpha
W_{\bm\alpha,\bm m}(q^\alpha)\right)$ by writing
$\partial_\alpha$ on $x^\alpha$ as multiplication by $\log x$.
\end{proof}

\begin{theorem}[Puiseux--log acceleration]
Suppose, as $x\downarrow0$,
\begin{equation}
 f(x)=a_0+
 \sum_{j=1}^J\sum_{s=0}^{m_j-1}
 a_{j,s}x^{\alpha_j}(\log x)^s
 +O\!\left(x^\beta(1+|\log x|^r)\right),
 \label{g1:eq:puiseux-log-expansion}
\end{equation}
where $\Re\alpha_j>0$ and $\beta>0$.  Then
\begin{equation}
 (\Rc_{\bm\alpha,\bm m}f)(x)
 =a_0+O\!\left(x^\beta(1+|\log x|^r)\right).
 \label{g1:eq:puiseux-acceleration}
\end{equation}
\end{theorem}

\begin{proof}
The finite sum of displayed modes is annihilated exactly.  For each fixed
$k\le N$, replacing $x$ by $q^kx$ preserves the stated remainder order, with a
constant depending on $q,k,\beta,r$.  Summing finitely many weighted remainders
proves the assertion.
\end{proof}

\begin{remark}[Relation to Richardson extrapolation]
This is geometric Richardson extrapolation in its natural spectral form.  The
usual integer-power Gaussian filter corresponds to
$\alpha_j=j$, $m_j=1$.  Repeated exponents handle logarithmic resonance without
introducing derivative data at the nodes.
\end{remark}

\begin{problem}[Stability-optimal modal design]
For a prescribed list of modes, \eqref{g1:eq:mode-polynomial} is algebraically
forced at minimum degree.  If extra degrees of freedom are allowed, determine
the polynomial $W$ with $W(1)=1$ and the required confluent zeros that minimizes
$\sum_k|w_k|$, or a weighted noise norm adapted to the expected precision of
$f(q^kx)$.  This is a finite convex optimization problem whose asymptotic
structure appears unexplored in the Fabius/Rvachev setting.
\end{problem}

\section{Stability, boundary layers, and regular variation}
\label{g1:sec:stability}

\subsection{Exact total variation}

Because the signs of \eqref{g1:eq:lambda} alternate, the total variation can be
read from the row polynomial at $-1$.

\begin{theorem}[Exact noise amplification]
\label{g1:thm:row-norm}
For $0<q<1$,
\begin{equation}
 K_n(q):=\sum_{k=0}^n|\lambda_{n,k}^{(q)}|
 =\frac{(-q;q)_n}{(q;q)_n}
 =\prod_{j=1}^n\frac{1+q^j}{1-q^j}.
 \label{g1:eq:K-n}
\end{equation}
Consequently
\begin{equation}
 K_n(q)\uparrow K_\infty(q):=\frac{(-q;q)_\infty}{(q;q)_\infty}<\infty.
 \label{g1:eq:K-infinity}
\end{equation}
If each sample has absolute error at most $\delta$, the extrapolated absolute
error is at most $K_n(q)\delta$.
\end{theorem}

\begin{proof}
Since
$\Sign\lambda_{n,k}=(-1)^{n-k}$,
\[
 \sum_k|\lambda_{n,k}|=(-1)^n\sum_k\lambda_{n,k}(-1)^k
 =(-1)^nW_{n,q}(-1).
\]
Substitute \eqref{g1:eq:row-polynomial}.  Monotone convergence of the positive
product gives \eqref{g1:eq:K-infinity}.
\end{proof}

For reference,
\[
 K_\infty(1/4)\approx1.9692603537,
 \qquad K_\infty(1/2)\approx8.2559879358,
 \qquad K_\infty(3/4)\approx803.0132718.
\]
Thus the Fourier base $Q=1/4$ associated with dyadic Rvachev scaling is much
better conditioned than a direct dyadic $q=1/2$ extrapolation.

\subsection{The singular regime $q\uparrow1$}

Two limits must be distinguished.  For fixed $n$,
\begin{equation}
 K_n(q)\sim\frac{2^n}{n!(1-q)^n}
 \qquad(q\uparrow1).
 \label{g1:eq:fixed-n-q1}
\end{equation}
If the order tends to infinity first, the growth is nonperturbative.

\begin{theorem}[Modular asymptotic]
Write $q=e^{-t}$ with $t\downarrow0$.  Then
\begin{equation}
 K_\infty(e^{-t})
 \sim\frac{\sqrt t}{2\sqrt\pi}
 \exp\!\left(\frac{\pi^2}{4t}\right).
 \label{g1:eq:K-modular}
\end{equation}
\end{theorem}

\begin{proof}
Use
$(-q;q)_\infty=(q^2;q^2)_\infty/(q;q)_\infty$, so
\[
 K_\infty(q)=\frac{(q^2;q^2)_\infty}{(q;q)_\infty^2}.
\]
The standard Dedekind-eta asymptotic
\[
 (e^{-t};e^{-t})_\infty
 \sim\sqrt{\frac{2\pi}{t}}
 \exp\!\left(-\frac{\pi^2}{6t}+\frac{t}{24}\right)
\]
applied at $t$ and $2t$ gives \eqref{g1:eq:K-modular}; the elementary exponential
corrections cancel.
\end{proof}

The plot in \cref{g1:fig:stability} shows the practical consequence: ratios close
to one are unsuitable for high-order accumulation-point extrapolation unless
precision grows very rapidly.

\subsection{The reversed-row boundary layer}

Put $r=n-k$.  For fixed $r$,
\begin{equation}
 \lambda_{n,n-r}^{(q)}\longrightarrow
 \beta_r^{(q)}:=
 \frac{(-1)^rq^{\binom{r+1}{2}}}{(q;q)_\infty(q;q)_r}.
 \label{g1:eq:beta}
\end{equation}
The limiting signed generating function is
\begin{equation}
 \sum_{r=0}^{\infty}\beta_r^{(q)}z^r
 =\frac{(qz;q)_\infty}{(q;q)_\infty}.
 \label{g1:eq:beta-generating}
\end{equation}
In particular, its mass is $1$ and its total variation is $K_\infty(q)$.

\begin{theorem}[$\ell^1$ boundary-layer convergence]
Extend the reversed row by zero for $r>n$.  Then
\begin{equation}
 \sum_{r=0}^{\infty}
 \left|\lambda_{n,n-r}^{(q)}-\beta_r^{(q)}\right|
 \longrightarrow0.
 \label{g1:eq:l1-boundary}
\end{equation}
\end{theorem}

\begin{proof}
Pointwise convergence follows from $(q;q)_{n-r}\to(q;q)_\infty$.  After
multiplication by $(-1)^r$, all terms are nonnegative.  Their total masses are
$K_n(q)$ and $K_\infty(q)$, respectively, and the former tends to the latter.
Scheffe's lemma for counting measure gives convergence in $\ell^1$.
\end{proof}

This theorem explains the stability mechanism.  For large $n$, the effective
weights sit a bounded number of steps behind the smallest node $Aq^n$; the
large-node coefficients are superexponentially small.

\subsection{Convergence for merely continuous data}

Analyticity is not needed for consistency.

\begin{theorem}[Continuous accumulation-point convergence]
If $f$ is continuous on $[0,A]$, then
\begin{equation}
 (\Gc_{n,q}f)(A)\longrightarrow f(0).
 \label{g1:eq:continuous-convergence}
\end{equation}
More quantitatively, if
$\omega_f(\delta)=\sup_{0\le x\le\delta}|f(x)-f(0)|$, then for
$1\le J\le n+1$,
\begin{align}
 |(\Gc_{n,q}f)(A)-f(0)|
 &\le K_\infty(q)\,\omega_f(Aq^J)\notag\\
 &\quad+\frac{2\|f\|_\infty J}{(q;q)_\infty^2}
 q^{\binom{n-J+2}{2}}.
 \label{g1:eq:modulus-bound}
\end{align}
\end{theorem}

\begin{proof}
Split the weighted sum into $k\ge J$ and $k<J$.  On the first part, every node
is at most $Aq^J$, so the total is bounded by $K_n\omega_f(Aq^J)$.  On the
second part,
\[
 |\lambda_{n,k}^{(q)}|
 \le\frac{q^{\binom{n-k+1}{2}}}{(q;q)_\infty^2}
 \le\frac{q^{\binom{n-J+2}{2}}}{(q;q)_\infty^2}.
\]
There are at most $J$ such terms, and
$|f(Aq^k)-f(0)|\le2\|f\|_\infty$.  First fix $J$ and let $n\to\infty$, then let
$J\to\infty$.
\end{proof}

\subsection{Regularly varying endpoint germs}

A function $f$ is regularly varying at $0+$ with index $\alpha\in\RealNumbers$ if
$f(tx)/f(t)\to x^\alpha$ as $t\downarrow0$ for every $x>0$, with $f$ eventually
of one sign.  Standard background is in \cite{BinghamGoldieTeugels}.

\begin{theorem}[Regular-variation transfer]
Let $f$ be regularly varying at $0+$ with index $\alpha$.  Then, for fixed
$A>0$,
\begin{equation}
 \frac{(\Gc_{n,q}f)(A)}{f(Aq^n)}
 \longrightarrow C_q(\alpha):=
 \frac{(q^{1-\alpha};q)_\infty}{(q;q)_\infty}.
 \label{g1:eq:regular-variation}
\end{equation}
For a positive integer $\alpha=m$, the limit is zero because the numerator
product contains the factor $1-q^0$.
\end{theorem}

\begin{proof}
Reverse the row:
\[
 \frac{(\Gc_{n,q}f)(A)}{f(Aq^n)}
 =\sum_{r=0}^n\lambda_{n,n-r}^{(q)}
   \frac{f(Aq^{n-r})}{f(Aq^n)}.
\]
For fixed $r$, regular variation gives a limit $q^{-\alpha r}$, while
\eqref{g1:eq:beta} gives the coefficient limit.  Potter bounds dominate the
ratio by $Cq^{-(\alpha+\varepsilon)r}$ away from a finite initial segment; the
factor $q^{r(r+1)/2}$ in $\beta_r$ makes the resulting series summable for every
$\alpha$.  The finitely many terms corresponding to nodes bounded away from
zero have superexponentially small coefficients.  Dominated convergence and
\eqref{g1:eq:beta-generating} now give
\[
 \sum_{r\ge0}\beta_r^{(q)}q^{-\alpha r}
 =\frac{(q^{1-\alpha};q)_\infty}{(q;q)_\infty}.
\]
\end{proof}

For pure powers, \eqref{g1:eq:regular-variation} is already visible in the exact
formula \eqref{g1:eq:pure-power}.  The theorem says that slowly varying factors do
not change the limiting multiplier.

\section{Smooth remainders and geometric-knot B-splines}
\label{g1:sec:bspline}

\subsection{Hermite--Genocchi representation}

The exact Taylor tail is best for analytic data.  For finite smoothness, a
divided-difference representation reveals positivity and sign.

\begin{theorem}[Geometric Hermite--Genocchi remainder]
Let $f\in C^{n+1}[0,A]$ and let
$(T_*,T_0,\ldots,T_n)$ be uniformly distributed on the $(n+1)$-simplex,
that is, Dirichlet with all parameters equal to $1$.  Put
\begin{equation}
 Y_{n,A,q}:=A\sum_{k=0}^nq^kT_k.
 \label{g1:eq:Y-dirichlet}
\end{equation}
Then
\begin{equation}
 (\Gc_{n,q}f)(A)-f(0)
 =\frac{(-1)^nA^{n+1}q^{\binom{n+1}{2}}}{(n+1)!}
 \Expectation\bigl[f^{(n+1)}(Y_{n,A,q})\bigr].
 \label{g1:eq:HG-remainder}
\end{equation}
\end{theorem}

\begin{proof}
Let $I_nf$ be the interpolating polynomial through $A,Aq,\ldots,Aq^n$.
The ordinary interpolation remainder at $0$ gives
\begin{align*}
 I_nf(0)-f(0)
 &=(-1)^nA^{n+1}q^{\binom{n+1}{2}}
 f[0,A,Aq,\ldots,Aq^n].
\end{align*}
The Hermite--Genocchi formula for a divided difference of order $n+1$ is
\[
 f[0,A,Aq,\ldots,Aq^n]
 =\frac1{(n+1)!}
 \Expectation f^{(n+1)}\!\left(A\sum_{k=0}^nq^kT_k\right).
\]
Since $I_nf(0)=(\Gc_{n,q}f)(A)$, the result follows.
\end{proof}

\begin{corollary}[Sharp derivative bound and sign]
Under the hypotheses of the theorem,
\begin{equation}
 |(\Gc_{n,q}f)(A)-f(0)|
 \le\frac{A^{n+1}q^{\binom{n+1}{2}}}{(n+1)!}
 \|f^{(n+1)}\|_{L^\infty(0,A)}.
 \label{g1:eq:smooth-bound}
\end{equation}
If $f^{(n+1)}$ has a fixed nonzero sign on $(0,A)$, then the remainder has the
sign of $(-1)^nf^{(n+1)}$.
\end{corollary}

\begin{proof}
The Hermite--Genocchi formula expresses the remainder as
$(-1)^nA^{n+1}q^{\binom{n+1}{2}}/(n+1)!$ times an average of
$f^{(n+1)}$ over points in $[0,A]$.  Taking absolute values gives the
bound, while a fixed nonzero sign of the derivative fixes the sign of
that average.
\end{proof}

\subsection{Peano kernel and a positive spline density}

For $n\ge1$, define the Peano kernel
\begin{equation}
 K_{n,A,q}(t):=\frac1{n!}\sum_{k=0}^n
 \lambda_{n,k}^{(q)}(Aq^k-t)_+^n,
 \qquad 0\le t\le A.
 \label{g1:eq:Peano-kernel}
\end{equation}
Then
\[
 (\Gc_{n,q}f)(A)-f(0)=\int_0^Af^{(n+1)}(t)K_{n,A,q}(t)\Differential t.
\]
Normalize it by
\begin{equation}
 B_{n,A,q}(t):=
 \frac{(-1)^n(n+1)!}{A^{n+1}q^{\binom{n+1}{2}}}K_{n,A,q}(t).
 \label{g1:eq:B-density}
\end{equation}

\begin{theorem}[Geometric-knot B-spline]
The function $B_{n,A,q}$ is the probability density of
$Y_{n,A,q}$ in \eqref{g1:eq:Y-dirichlet}.  In particular,
\begin{enumerate}[label=(\roman*),leftmargin=2em]
\item $B_{n,A,q}\ge0$ and is supported on $[0,A]$;
\item $\int_0^AB_{n,A,q}(t)\Differential t=1$;
\item it is the classical univariate B-spline associated with the knots
$0,Aq^n,\ldots,Aq,A$, up to the displayed normalization;
\item it is log-concave and hence unimodal on the interior of its support.
\end{enumerate}
\end{theorem}

\begin{proof}
Comparing the Peano representation with \eqref{g1:eq:HG-remainder} for arbitrary
continuous $f^{(n+1)}$ identifies the normalized kernel with the law of
$Y_{n,A,q}$.  This is also the standard Curry--Schoenberg interpretation of a
B-spline as the density of a Dirichlet average of its knots; see
\cite{deBoor}.  Positivity, support, and normalization follow immediately.
The uniform density on a simplex is log-concave, and linear images preserve
log-concavity by the Prekopa theorem, proving the last assertion.
\end{proof}

\begin{remark}[Not the same as a ``quantum B-spline'']
The density above is a classical B-spline whose knots happen to be geometric.
It should be distinguished from the quantum B-splines of
Simeonov--Goldman \cite{SimeonovGoldman}, where continuity conditions are
formulated using $q$-derivatives.  The two theories meet through divided
differences and Jackson calculus, but their definitions are different.  The
geometric-knot density is the one canonically forced by the interpolation
remainder in this report.
\end{remark}

\subsection{Exact moments}

\begin{theorem}[Gaussian-binomial B-spline moments]
For every $r\in\NaturalNumbers$,
\begin{equation}
 \Expectation Y_{n,A,q}^r
 =\frac{r!(n+1)!}{(n+r+1)!}
 A^r\GaussianBinomial{n+r}{n}{q}.
 \label{g1:eq:Y-moments}
\end{equation}
Equivalently,
\begin{equation}
 \int_0^At^rB_{n,A,q}(t)\Differential t
 =\frac{r!(n+1)!}{(n+r+1)!}
 A^r\GaussianBinomial{n+r}{n}{q}.
\end{equation}
\end{theorem}

\begin{proof}
For a Dirichlet$(1,\ldots,1)$ vector with $n+2$ components,
\[
 \Expectation\left(\sum_i x_iT_i\right)^r
 =\frac{r!(n+1)!}{(n+r+1)!}h_r(x_0,\ldots,x_{n+1}).
\]
Take the knots $0,A,Aq,\ldots,Aq^n$ and apply
\eqref{gq:eq:h-principal}.
\end{proof}

Thus the Gaussian coefficient is not only an interpolation residual: it is the
moment sequence of the positive Peano-kernel distribution after a simple beta
normalization.

\subsection{A $q$-Pochhammer scaling limit}

The geometric B-spline has a nontrivial boundary-scale limit.  Let
$E_*,E_0,E_1,\ldots$ be independent exponential random variables of mean $1$.

\begin{theorem}[Scaling limit of the geometric B-spline]
For fixed $0<q<1$,
\begin{equation}
 \frac{n+2}{A}Y_{n,A,q}
 \xrightarrow{\ d\ }
 Z_q:=\sum_{k=0}^{\infty}q^kE_k.
 \label{g1:eq:Y-scaling-limit}
\end{equation}
The limit has moments and Laplace transform
\begin{align}
 \Expectation Z_q^r&=\frac{r!}{(q;q)_r},
 \label{g1:eq:Z-moments}\\
 \Expectation e^{-sZ_q}&=\frac1{(-s;q)_\infty},\qquad s\ge0.
 \label{g1:eq:Z-Laplace}
\end{align}
For $y>0$ its density admits the absolutely convergent expansion
\begin{equation}
 p_q(y)=\frac1{(q;q)_\infty}
 \sum_{k=0}^{\infty}
 \frac{(-1)^kq^{\binom{k}{2}}}{(q;q)_k}
 e^{-q^{-k}y}.
 \label{g1:eq:Z-density}
\end{equation}
\end{theorem}

\begin{proof}
Use the exponential representation of a uniform Dirichlet vector:
\[
 (T_*,T_0,\ldots,T_n)
 \overset d=
 \frac{(E_*,E_0,\ldots,E_n)}{E_*+E_0+\cdots+E_n}.
\]
The denominator divided by $n+2$ converges almost surely to $1$, while
$\sum_{k=0}^nq^kE_k$ converges almost surely and in every $L^r$ to $Z_q$.
Slutsky's theorem proves \eqref{g1:eq:Y-scaling-limit}.  Independence gives
\[
 \Expectation e^{-sZ_q}=\prod_{k\ge0}(1+sq^k)^{-1}=(-s;q)_\infty^{-1}.
\]
Expanding $1/(t;q)_\infty=\sum_{r\ge0}t^r/(q;q)_r$ around $t=0$ and comparing
moment-generating coefficients gives \eqref{g1:eq:Z-moments}.

For the density, first decompose the finite product
$\prod_{j=0}^{N}(1+sq^j)^{-1}$ into simple fractions.  Its coefficient at the
pole $s=-q^{-k}$ is the finite version of
\[
 \frac1{(-s;q)_\infty}
 =\sum_{k\ge0}
 \frac{(-1)^kq^{\binom{k+1}{2}}}
 {(q;q)_k(q;q)_\infty}\frac1{1+sq^k}.
\]
For $s\ge0$, the finite coefficients converge to the displayed coefficients
and are dominated by a constant times
$q^{\binom{k+1}{2}}/(q;q)_k$; dominated convergence therefore gives the
infinite partial-fraction identity.  Termwise Laplace inversion gives
\eqref{g1:eq:Z-density}.  For each $y>0$, the Gaussian factor in $k$ together with
$e^{-q^{-k}y}$ gives absolute convergence and justifies the inversion.
\end{proof}

\begin{remark}
The limit satisfies the perpetuity equation
\[
 Z_q\overset d=E_0+qZ_q',
\]
where $Z_q'$ is an independent copy.  The $q$-Pochhammer product is therefore
both an interpolation object and the Laplace transform of the boundary-scale
limit of the associated B-splines.
\end{remark}

\section{Geometric-uniform random sums and finite splines}
\label{g1:sec:uniform-sums}

\subsection{The generic law}

Let $0<b<1$ and let $U_0,U_1,\ldots$ be independent uniforms on $[0,1]$.
Define
\begin{equation}
 X_b:=(1-b)\sum_{j=0}^{\infty}b^jU_j.
 \label{g1:eq:Xb}
\end{equation}
The weights sum to $1$, so $X_b\in[0,1]$.  Splitting off the first summand gives
the distributional self-similarity
\begin{equation}
 X_b\overset d=(1-b)U_0+bX_b',
 \label{g1:eq:Xb-selfsimilar}
\end{equation}
where $X_b'$ is an independent copy.  If $F_b$ is the distribution function,
then
\begin{equation}
 F_b(x)=\int_0^1F_b\!\left(\frac{x-(1-b)u}{b}\right)\Differential u,
 \label{g1:eq:Fb-integral}
\end{equation}
with the conventions $F_b=0$ to the left of $0$ and $F_b=1$ to the right of
$1$.  At $b=1/2$, this is the Fabius distribution.

\subsection{Finite generalized Irwin--Hall splines}

Let
\[
 X_{b,N}:=(1-b)\sum_{j=0}^{N-1}b^jU_j,
 \qquad a_j=(1-b)b^j.
\]
This is a convolution of uniforms on intervals of geometrically decreasing
width.

\begin{theorem}[Finite box-spline formula]
For $N\ge1$, the density and distribution function of $X_{b,N}$ are
\begin{align}
 f_{b,N}(x)
 &=\frac1{(N-1)!\prod_{j=0}^{N-1}a_j}
 \sum_{\varepsilon\in\{0,1\}^N}
 (-1)^{|\varepsilon|}
 \left(x-\sum_{j=0}^{N-1}\varepsilon_ja_j\right)_+^{N-1},
 \label{g1:eq:finite-density}\\
 F_{b,N}(x)
 &=\frac1{N!\prod_{j=0}^{N-1}a_j}
 \sum_{\varepsilon\in\{0,1\}^N}
 (-1)^{|\varepsilon|}
 \left(x-\sum_{j=0}^{N-1}\varepsilon_ja_j\right)_+^N,
 \label{g1:eq:finite-CDF}
\end{align}
where
\begin{equation}
 \prod_{j=0}^{N-1}a_j=(1-b)^Nb^{N(N-1)/2}.
 \label{g1:eq:product-widths}
\end{equation}
\end{theorem}

\begin{proof}
The density of a uniform on $[0,a]$ is
$a^{-1}(H(x)-H(x-a))$, where $H$ is the Heaviside function.  Convolving $N$
such factors and expanding the product of translation differences gives the
alternating subset sum.  The $N$-fold convolution of $H$ is
$x_+^{N-1}/(N-1)!$.  Integration gives the CDF, and
\eqref{g1:eq:product-widths} is immediate.
\end{proof}

For a general $b$, the knots are subset sums of a geometric progression.  At
the dyadic value they collapse to an arithmetic comb, and the coefficient sign
is the Thue--Morse sign.

\begin{corollary}[Dyadic Thue--Morse collapse]
For $b=1/2$,
\begin{equation}
 F_{1/2,N}(x)
 =\frac{2^{N(N+1)/2}}{N!}
 \sum_{r=0}^{2^N-1}(-1)^{\BinaryDigitSum(r)}
 \left(x-\frac{r}{2^N}\right)_+^N,
 \label{g1:eq:dyadic-finite-CDF}
\end{equation}
where $\BinaryDigitSum(r)$ is the binary digit sum of $r$.
\end{corollary}

\begin{proof}
Here $a_j=2^{-(j+1)}$.  Every subset sum is uniquely $r/2^N$, and the subset
cardinality has the parity of $\BinaryDigitSum(r)$.  The product of widths is
$2^{-N(N+1)/2}$.
\end{proof}

This identity explains why a geometric collection of widths produces an
arithmetic dyadic knot comb: binary representation converts subset selection
into position on the uniform grid, while the inclusion--exclusion sign becomes
the Thue--Morse sequence.

\subsection{Convergence to the infinite law}

The omitted tail is bounded deterministically:
\[
 0\le X_b-X_{b,N}\le b^N.
\]
Hence
\begin{equation}
 F_b(x)\le F_{b,N}(x)\le F_b(x+b^N).
 \label{g1:eq:CDF-sandwich}
\end{equation}
Because $F_b$ is continuous, $F_{b,N}\to F_b$ uniformly.  Thus the Fabius
function is a uniform limit of explicitly known compactly supported splines
whose coefficients are a Thue--Morse finite difference.

\subsection{Centered law, Bernoulli cumulants, and sinc products}

Put
\begin{equation}
 Z_b:=2X_b-1
 =2(1-b)\sum_{j=0}^{\infty}b^j(U_j-1/2).
 \label{g1:eq:Zb}
\end{equation}
Then $Z_b$ is symmetric and supported on $[-1,1]$.  Under the Fourier
convention
\[
 \FourierTwoPi{u_b}(z)=\Expectation\,\EulerE^{-2\pi\ImaginaryUnit zZ_b},
\]
its characteristic function is
\begin{equation}
 \FourierTwoPi{u_b}(z)
 =\prod_{j=0}^{\infty}
 \SincRad\!\bigl(2\pi(1-b)b^jz\bigr),
 \qquad \SincRad w:=\frac{\sin w}{w}.
 \label{g1:eq:generic-sinc-product}
\end{equation}
At $b=1/2$, this is
\begin{equation}
 \FourierTwoPi{\RvachevUp}(z)=\prod_{j=0}^{\infty}\SincRad(\pi z/2^j),
 \label{g1:eq:up-sinc}
\end{equation}
which is the standard Rvachev product.

The centered uniform on $[-1/2,1/2]$ has cumulants
$\kappa_m=B_m/m$ for $m\ge2$, where $B_m$ are Bernoulli numbers and odd
cumulants vanish.  Independence and scaling give the following explicit
$q$-denominators.

\begin{proposition}[Cumulants of the geometric-uniform law]
For $m\ge1$,
\begin{equation}
 \kappa_{2m}(Z_b)
 =\frac{[2(1-b)]^{2m}B_{2m}}
 {2m(1-b^{2m})},
 \qquad \kappa_{2m+1}(Z_b)=0.
 \label{g1:eq:Zb-cumulants}
\end{equation}
Consequently every even moment is an explicit complete Bell polynomial in the
numbers \eqref{g1:eq:Zb-cumulants}.
\end{proposition}

\begin{proof}
Cumulants add under independent sums and scale homogeneously.  Therefore
\[
 \kappa_{2m}(Z_b)
 =[2(1-b)]^{2m}\frac{B_{2m}}{2m}
 \sum_{j\ge0}b^{2mj},
\]
and the geometric sum gives \eqref{g1:eq:Zb-cumulants}.
\end{proof}

The factors $1-b^{2m}$ are the simplest $q$-analog denominators.  In Fourier
space the same even degree $2m$ is responsible for the squared interpolation
base developed in \cref{g1:sec:fourier}.

\section{Fabius and Rvachev functions at the flat boundary}
\label{g1:sec:fabius}

\subsection{Definitions and endpoint differential scaling}

Let $F$ be the Fabius distribution function on $[0,1]$.  It is the CDF of
$X_{1/2}$ and satisfies
\begin{equation}
 F(0)=0,\qquad F(1)=1,\qquad F(1-x)=1-F(x),
 \label{g1:eq:F-symmetry}
\end{equation}
and, for $0\le x\le1/2$,
\begin{equation}
 F'(x)=2F(2x).
 \label{g1:eq:F-DE}
\end{equation}
The compactly supported Rvachev function is related by
\begin{equation}
 \RvachevUp(x)=F(1-|x|),\qquad |x|\le1,
 \label{g1:eq:Up-F}
\end{equation}
with $\RvachevUp(x)=0$ for $|x|\ge1$.  These normalizations agree with the primary
Fabius/Rvachev literature \cite{Fabius1966,Rvachev1990,AriasDeReyna} and with
the inspected ProveIt exposition \cite{ProveItFabius}.

Repeated differentiation of \eqref{g1:eq:F-DE} gives the exact endpoint scaling.

\begin{lemma}[Endpoint derivatives]
For $m\ge1$ and $0\le x\le2^{-m}$,
\begin{equation}
 F^{(m)}(x)=2^{m(m+1)/2}F(2^mx).
 \label{g1:eq:F-derivatives}
\end{equation}
\end{lemma}

\begin{proof}
The case $m=1$ is \eqref{g1:eq:F-DE}.  If the result holds for $m$, then on
$x\le2^{-(m+1)}$,
\[
 F^{(m+1)}(x)=2^{m(m+1)/2}2^mF'(2^mx)
 =2^{m(m+1)/2+m+1}F(2^{m+1}x),
\]
which has exponent $(m+1)(m+2)/2$.
\end{proof}

\subsection{An exact boundary-comb identity}

The smooth Peano remainder becomes self-similar when $f=F$.

\begin{theorem}[Fabius boundary filter]
Let $n\ge0$, $0<A\le2^{-(n+1)}$, and let
$(T_*,T_0,\ldots,T_n)$ be uniform Dirichlet.  Put
\begin{equation}
 S_n:=\sum_{k=0}^n2^{-k}T_k\in[0,1].
 \label{g1:eq:Sn}
\end{equation}
Then
\begin{equation}
 \sum_{k=0}^n\lambda_{n,k}^{(1/2)}F(A2^{-k})
 =\frac{(-1)^n(2A)^{n+1}}{(n+1)!}
 \Expectation F(2^{n+1}AS_n).
 \label{g1:eq:F-boundary-filter}
\end{equation}
In particular, the left side has strict sign $(-1)^n$ and obeys
\begin{equation}
 0<(-1)^n\sum_{k=0}^n\lambda_{n,k}^{(1/2)}F(A2^{-k})
 \le\frac{(2A)^{n+1}}{(n+1)!}.
 \label{g1:eq:F-boundary-bound}
\end{equation}
\end{theorem}

\begin{proof}
Apply \eqref{g1:eq:HG-remainder} with $q=1/2$ and $f=F$.  Since $F(0)=0$,
\[
 \sum_k\lambda_{n,k}^{(1/2)}F(A2^{-k})
 =\frac{(-1)^nA^{n+1}2^{-n(n+1)/2}}{(n+1)!}
 \Expectation F^{(n+1)}(AS_n).
\]
The hypothesis on $A$ allows \eqref{g1:eq:F-derivatives} with $m=n+1$ throughout
the support of $AS_n$.  The powers of $2$ reduce to $2^{n+1}$, giving
\eqref{g1:eq:F-boundary-filter}.  The random argument lies in $[0,1]$, so
$0\le F\le1$.  It is positive almost surely because $S_n>0$ almost surely,
which gives strict sign and the bound.
\end{proof}

\begin{corollary}[Exact dyadic form]
If $M\ge n+1$, then
\begin{equation}
 \sum_{k=0}^n\lambda_{n,k}^{(1/2)}F(2^{-(M+k)})
 =\frac{(-1)^n2^{-(M-1)(n+1)}}{(n+1)!}
 \Expectation F(2^{n+1-M}S_n).
 \label{g1:eq:F-dyadic-filter}
\end{equation}
For $M=n+1$, the ratio of the absolute value to the upper bound is exactly
$\Expectation F(S_n)$.
\end{corollary}

\begin{proof}
Set $A=2^{-M}$ in \cref{g1:eq:F-boundary-filter}; the hypothesis
$M\ge n+1$ is exactly the required boundary-range condition.  Then
$(2A)^{n+1}=2^{-(M-1)(n+1)}$, proving the formula.  At $M=n+1$ the
argument of $F$ is $S_n$, and division by the preceding theorem's upper
bound leaves precisely $\Expectation F(S_n)$.
\end{proof}

Through \eqref{g1:eq:Up-F}, the same result is an identity on the left boundary of
$\RvachevUp$:
\begin{equation}
 \sum_{k=0}^n\lambda_{n,k}^{(1/2)}
 \RvachevUp(-1+A2^{-k})
 =\frac{(-1)^n(2A)^{n+1}}{(n+1)!}
 \Expectation F(2^{n+1}AS_n).
 \label{g1:eq:Up-boundary-filter}
\end{equation}
Thus a signed combination of values of a nowhere-analytic $C^\infty$ function
is represented by a positive expectation and has a completely determined
sign.

\subsection{The boundary B-spline and its $q$-exponential scale}

The variable $S_n$ is $Y_{n,1,1/2}$ from \cref{g1:sec:bspline}.  Therefore
\begin{equation}
 (n+2)S_n\xrightarrow{d}Z_{1/2}
 =\sum_{k\ge0}2^{-k}E_k,
 \qquad
 \Expectation e^{-sZ_{1/2}}=\frac1{(-s;1/2)_\infty}.
 \label{g1:eq:Sn-limit}
\end{equation}
This identifies the random scale sampled by the normalized Fabius remainder.
The typical size of $S_n$ is $2/(n+2)$, but its limiting rescaled distribution
is not degenerate; it is the $q$-Pochhammer perpetuity in
\eqref{g1:eq:Z-Laplace}.

\subsection{Exact rational computation}

All dyadic values of $F$ are rational.  A convenient triangular recurrence,
also used in the linked monograph, is
\begin{equation}
 V_0=1,\qquad
 V_m:=F(2^{-m})
 =\frac{2^{-\binom m2}}{2^m-1}
 \sum_{k=0}^{m-1}
 \frac{2^{\binom k2}}{(m-k+1)!}V_k.
 \label{g1:eq:V-recurrence}
\end{equation}
It begins
\[
 V_0=1,\quad V_1=\frac12,\quad V_2=\frac5{72},\quad
 V_3=\frac1{288},\quad V_4=\frac{143}{2073600}.
\]
Because both \eqref{g1:eq:V-recurrence} and the dyadic weights
$\lambda_{n,k}^{(1/2)}$ are rational, \eqref{g1:eq:F-dyadic-filter} can be checked
with exact integer arithmetic.  The companion Python program does so before
switching to multiprecision arithmetic for larger plotted orders.

\section{Fourier-scale interpolation for the Rvachev function}
\label{g1:sec:fourier}

\subsection{The base-squaring principle}

Let $\Phi$ be an even entire function normalized by $\Phi(0)=1$, and write
\begin{equation}
 \Phi(z)=\sum_{m=0}^{\infty}c_mz^{2m}.
 \label{g1:eq:Phi-even}
\end{equation}
Fix a scale ratio $0<c<1$ and set
\begin{equation}
 Q:=c^2,\qquad H_z(t):=\Phi(z\sqrt t)
 =\sum_{m=0}^{\infty}c_mz^{2m}t^m.
 \label{g1:eq:Hzt}
\end{equation}
Evenness makes $H_z$ entire in $t$, independent of the choice of square root.
Moreover
\begin{equation}
 H_z(Q^k)=\Phi(c^kz).
 \label{g1:eq:scale-samples}
\end{equation}
Therefore geometric interpolation of Fourier scales $c^kz$ is interpolation on
the comb $1,Q,Q^2,\ldots$, not on base $c$ itself.  In the dyadic Rvachev
problem, $c=1/2$ and the interpolation base is $Q=1/4$.

\subsection{Exact moment-flat filter}

Let $\lambda_{n,k}^{(Q)}$ denote the origin weights with base $Q$.

\begin{theorem}[Even-entire scale filter]
For $H_z$ as in \eqref{g1:eq:Hzt},
\begin{align}
 \sum_{k=0}^n\lambda_{n,k}^{(Q)}\Phi(c^kz)
 &=1+(-1)^nQ^{\binom{n+1}{2}}
 \sum_{r=0}^{\infty}\GaussianBinomial{n+r}{n}{Q}
 c_{n+1+r}z^{2n+2+2r}.
 \label{g1:eq:even-entire-filter}
\end{align}
In particular, the residual has a zero of order at least $2n+2$ at $z=0$.
\end{theorem}

\begin{proof}
Apply \eqref{g1:eq:exact-analytic-tail} to the entire function $H_z(t)$ at
$A=1$ and base $Q$.
\end{proof}

For a characteristic function under the convention
$\Phi(z)=\Expectation e^{-2\pi izX}$ with a symmetric random variable $X$,
\begin{equation}
 c_m=(-1)^m\frac{(2\pi)^{2m}\mu_{2m}}{(2m)!},
 \qquad \mu_{2m}=\Expectation X^{2m}.
 \label{g1:eq:Fourier-coeff-moments}
\end{equation}
Hence the first nonzero term of the filtered residual is always
\begin{equation}
 -Q^{\binom{n+1}{2}}
 \frac{(2\pi z)^{2n+2}\mu_{2n+2}}{(2n+2)!}.
 \label{g1:eq:leading-Fourier-residual}
\end{equation}
For real $z$ near zero this is negative, regardless of the parity of $n$.

\subsection{Application to $\RvachevUp$}

For the Rvachev function,
\[
 \Phi(z)=\FourierTwoPi{\RvachevUp}(z)
 =\prod_{j=0}^{\infty}\SincRad(\pi z/2^j).
\]
Taking $c=1/2$ and $Q=1/4$ in
\eqref{g1:eq:even-entire-filter} gives an exact finite multiscale identity.  Its
uniform sample-noise amplification is bounded by
$K_\infty(1/4)\approx1.96926$, which is notably mild.

More generally, the centered geometric-uniform density $u_b$ in
\eqref{g1:eq:generic-sinc-product} has an even transform, and samples at scales
$b^kz$ are governed by the interpolation base $Q=b^2$.  The cumulants
\eqref{g1:eq:Zb-cumulants} determine the coefficients through Bell polynomials,
so the filter simultaneously links Gaussian coefficients, Bernoulli numbers,
Bell polynomials, and the infinite sinc product.

\subsection{Entire Newton reconstruction of all scales}

The finite filter recovers the constant term approximately.  The infinite
Newton series reconstructs the complete scale function.

\begin{theorem}[Scale reconstruction from a geometric sample ray]
Let $\Phi$ be even entire, fix $z\in\ComplexNumbers$, and set $H_z$ by
\eqref{g1:eq:Hzt}.  Then for every $t\in\ComplexNumbers$,
\begin{equation}
 \Phi(z\sqrt t)
 =\sum_{j=0}^{\infty}
 \frac{D_Q^jH_z(1)}{[j]_Q!}
 \prod_{r=0}^{j-1}(t-Q^r),
 \label{g1:eq:Phi-scale-Newton}
\end{equation}
locally uniformly in $t$.  Each coefficient is a finite linear combination of
$\Phi(z),\Phi(cz),\ldots,\Phi(c^jz)$:
\begin{equation}
 \frac{D_Q^jH_z(1)}{[j]_Q!}
 =\sum_{k=0}^j
 \frac{(-1)^kQ^{-kj+\binom{k+1}{2}}}
 {(Q;Q)_k(Q;Q)_{j-k}}\Phi(c^kz).
 \label{g1:eq:Phi-scale-coeff}
\end{equation}
At $t=0$ this gives the exact infinite sum rule
\begin{equation}
 1=\sum_{j=0}^{\infty}
 \frac{(-1)^jQ^{\binom j2}}{[j]_Q!}D_Q^jH_z(1).
 \label{g1:eq:Phi-scale-origin}
\end{equation}
\end{theorem}

\begin{proof}
Apply the locally uniform Newton series
\eqref{gq:eq:infinite-newton-series}, with coefficients given by
\eqref{gq:eq:geometric-divided-difference}, to the entire function $H_z$
at $A=1$ and base $Q$.  Evaluating that series at $t=0$ gives the final
identity.
\end{proof}

This is a strong determinacy statement: at any fixed frequency $z$, the
multiscale values on one geometric ray determine the transform at every complex
scale parameter $t$.  For the Rvachev product, this turns the self-similar
sample sequence into a complete Newton coordinate system.

\section{Algorithms and numerical implementation}
\label{g1:sec:algorithms}

The formulas admit several computational realizations.  They should not be
used interchangeably without regard to conditioning.

\subsection{Barycentric interpolation away from zero}

For an arbitrary target $x$, use the second barycentric formula
\eqref{gq:eq:second-barycentric}.  The raw weights
$b_{n,k}^{A,q}=1/d_{n,k}^{A,q}$ can span a large range, so in floating-point
arithmetic they should be rescaled by a common factor.  Only their ratios enter
the formula.  For large $n$, compute logarithms of absolute values from
\eqref{gq:eq:barycentric-weight} and attach the sign $(-1)^k$.

At a node, return the corresponding datum directly rather than evaluating a
removable $0/0$ expression.

\subsection{Newton--Horner evaluation}

Once the coefficients
$c_k=D_q^kf(A)/[k]_q!$ are known, evaluate
\eqref{gq:eq:q-taylor-interpolant} in nested form:
\begin{equation}
 \Ic_nf(x)=c_0+(x-A)\bigl(c_1+(x-Aq)
 (c_2+\cdots+(x-Aq^{n-1})c_n)\cdots\bigr).
 \label{g1:eq:Newton-Horner}
\end{equation}
This uses $O(n)$ arithmetic per target.  Coefficients can be generated from
samples by a divided-difference table or by
\eqref{gq:eq:geometric-divided-difference}.

\subsection{Origin extrapolation by factor recursion}

For the special target $0$, the recursion \eqref{g1:eq:triangular-filter} is
preferable to solving for an interpolating polynomial.  In pseudocode:

\begin{algorithm}[Factorized geometric-comb filter]
Given $y_k=f(Aq^k)$ for $0\le k\le n$:
\begin{enumerate}[leftmargin=2em]
\item Set \texttt{work[k] = y[k]}.
\item For $j=1,\ldots,n$, replace
\[
 \texttt{work[k]}
 \leftarrow
 \frac{\texttt{work[k+1]}-q^j\texttt{work[k]}}{1-q^j},
 \qquad 0\le k\le n-j.
\]
\item Return \texttt{work[0]}.
\end{enumerate}
\end{algorithm}

The induced $\ell^\infty$ error factor of the $j$th stage is at most
$(1+q^j)/(1-q^j)$; their product is exactly $K_n(q)$.  Thus the recursive
algorithm exposes the same stability constant as the explicit row.

\subsection{Exact dyadic arithmetic}

When $q=1/2$ and all samples are rational, every weight is rational.  For the
Fabius sequence, \eqref{g1:eq:V-recurrence} and the filter can therefore be
implemented with arbitrary-precision integer fractions.  This is valuable for
separating true cancellation from roundoff.  Numerators and denominators grow
quickly, so the companion code uses exact fractions only for moderate order and
switches to high-precision \texttt{mpmath} arithmetic for plots.

\subsection{Precision budgeting}

Suppose the desired absolute output accuracy is $\varepsilon$ and each sample
is evaluated with absolute accuracy $\delta$.  A sufficient condition is
\begin{equation}
 \delta\le\frac{\varepsilon}{K_n(q)}.
 \label{g1:eq:precision-budget}
\end{equation}
For fixed $q<1$, the right side remains comparable to $\varepsilon$ as
$n\to\infty$.  Near $q=1$, \eqref{g1:eq:K-modular} shows that the required number
of guard digits grows like
\[
 \frac{\pi^2}{4(-\log q)\log 10}
\]
in the infinite-order regime.  Relative precision is subtler: if the samples
themselves are superflat, as for $F(2^{-m})$, a small absolute error may still
be a large relative error.  Exact rational or ball arithmetic is preferable in
that case.

\section{Numerical experiments}
\label{g1:sec:numerics}

The file \texttt{geometric\_comb\_experiments.py} included in the archive
reproduces all data and figures in this section.  It uses 140 decimal digits for
alternating sums and exact rational arithmetic for the first dyadic Fabius
checks.  The numerical experiments are confirmations of proved identities,
not substitutes for proofs.

\subsection{Moment cancellation and factor recursion}

For $q=1/2$ and orders $1\le n\le20$, the script verifies
\[
 \sum_{k=0}^n\lambda_{n,k}^{(1/2)}2^{-kd}=0,
 \qquad 1\le d\le n,
\]
and compares the first residual with
$(-1)^n2^{-n(n+1)/2}$.  The largest observed canceled moment is at the level of
$10^{-140}$, consistent with the working precision.  The direct weighted sum
and the triangular factor recursion agree to the same scale on the
nonpolynomial test function $e^x$.

\subsection{Noise amplification}

\begin{figure}[htbp]
\centering
\includegraphics[width=0.88\textwidth]{stability_growth.png}
\caption{The infinite-order total variation
$K_\infty(q)=(-q;q)_\infty/(q;q)_\infty$.  The vertical axis is logarithmic in
base $10$.  Stability is excellent for small and moderate $q$ but deteriorates
nonperturbatively as $q\uparrow1$.}
\label{g1:fig:stability}
\end{figure}

The data file also compares $K_\infty(e^{-t})$ to the asymptotic
\eqref{g1:eq:K-modular}.  For the plotted near-one values the relative difference
is already below the displayed precision, reflecting the exponentially small
modular correction.

\subsection{Regular-variation constants}

\begin{figure}[htbp]
\centering
\includegraphics[width=0.88\textwidth]{regular_variation_convergence.png}
\caption{For pure powers $x^\alpha$ at $q=1/2$, the normalized exact multiplier
$(q^{1-\alpha};q)_n/(q;q)_n$ converges geometrically to
$C_q(\alpha)$.  The three exponents illustrate positive and negative limiting
constants.}
\label{g1:fig:regular-variation}
\end{figure}

For example,
\begin{align*}
 C_{1/2}(0.7)&\approx 0.2493175488440439431,\\
 C_{1/2}(1.3)&\approx-0.1538705960439442305,\\
 C_{1/2}(2.4)&\approx 0.2913748265379615867.
\end{align*}
At positive integer exponents the limit is exactly zero, in agreement with
polynomial cancellation.

\subsection{Exact Fabius boundary remainders}

For $A=2^{-(n+1)}$, define
\[
 R_n:=\sum_{k=0}^n\lambda_{n,k}^{(1/2)}F(2^{-(n+1+k)}),
 \qquad
 B_n:=\frac{2^{-n(n+1)}}{(n+1)!}.
\]
The theorem gives $(-1)^nR_n>0$ and $|R_n|/B_n=\Expectation F(S_n)\in(0,1)$.  Exact
fraction arithmetic produces the following decimal summary.

\begin{table}[htbp]
\centering
\caption{Normalized exact dyadic Fabius boundary remainders.}
\label{g1:tab:Fabius-remainders}
\begin{tabular}{rcc}
\toprule
$n$ & sign of $R_n$ & $|R_n|/B_n$\\
\midrule
1  & $-$ & $0.500000000000$\\
2  & $+$ & $0.392067901235$\\
3  & $-$ & $0.285827999336$\\
4  & $+$ & $0.204528650771$\\
5  & $-$ & $0.146778506870$\\
6  & $+$ & $0.106554660592$\\
7  & $-$ & $0.078481279976$\\
8  & $+$ & $0.058671832771$\\
9  & $-$ & $0.044492131277$\\
10 & $+$ & $0.034188031436$\\
\bottomrule
\end{tabular}
\end{table}

\begin{figure}[htbp]
\centering
\includegraphics[width=0.88\textwidth]{fabius_boundary_ratio.png}
\caption{The normalized Fabius boundary remainder
$|R_n|/B_n=\Expectation F(S_n)$ through order $36$.  The decrease is much slower than
the already extracted factorial-supergeometric factor $B_n$; its finer
asymptotics are posed as a conjecture in \cref{g1:sec:frontiers}.}
\label{g1:fig:Fabius-ratio}
\end{figure}

\subsection{Reproducibility files}

The script writes four CSV files:
\begin{center}
\begin{tabular}{ll}
\toprule
File & Contents\\
\midrule
\texttt{moment\_checks.csv} & canceled moments and recursion comparison\\
\texttt{stability\_growth.csv} & $K_\infty$, modular approximation, ratio\\
\texttt{regular\_variation.csv} & finite multipliers and limiting constants\\
\texttt{fabius\_boundary.csv} & exact rational filters, bounds, sign checks\\
\bottomrule
\end{tabular}
\end{center}
The archive includes these files and the three PNG figures, so the LaTeX source
can be recompiled without rerunning Python.

\section{Research frontiers and conjectures}
\label{g1:sec:frontiers}

\subsection{A local limit theorem for geometric B-splines}

\cref{g1:eq:Y-scaling-limit} proves weak convergence.  The explicit spline and
limit densities suggest a substantially stronger statement.

\begin{conjecture}[Local density limit]
For every fixed $0<q<1$ and compact interval $K\subset(0,\infty)$,
\begin{equation}
 \sup_{y\in K}
 \left|
 \frac{A}{n+2}B_{n,A,q}\!\left(\frac{Ay}{n+2}\right)-p_q(y)
 \right|\longrightarrow0,
 \label{g1:eq:local-density-conjecture}
\end{equation}
where $p_q$ is given by \eqref{g1:eq:Z-density}.  The convergence should extend to
all derivatives on compact subsets of $(0,\infty)$ after the natural scaling.
\end{conjecture}

A proof might combine the exponential/Dirichlet representation with Fourier or
Laplace inversion and uniform control of the denominator
$(E_*+\cdots+E_n)/(n+2)$.  Such a theorem would turn finite interpolation
kernels into quantitative approximations of a canonical $q$-exponential
density.

\subsection{Fabius--Dirichlet small-value asymptotics}

Let
\[
 \rho_n:=\Expectation F(S_n)=\frac{|R_n|}{B_n}.
\]
The scaling limit $(n+2)S_n\Rightarrow Z_{1/2}$ suggests that $S_n$ is of order
$1/n$, while the Fabius function has a logarithmic-square flatness scale at
zero.  Random fluctuations of $Z_{1/2}$ should affect lower-order terms but not
the leading quadratic logarithm.

\begin{conjecture}[Leading decay of the normalized Fabius remainder]
\begin{equation}
 \log\rho_n
 \sim-\frac{(\log n)^2}{2\log2}
 \qquad(n\to\infty).
 \label{g1:eq:Fabius-rho-conjecture}
\end{equation}
More precisely, the next terms should be obtainable by combining the sharp
Lambert-$W$ endpoint expansion of $F$ developed elsewhere in the ProveIt
project with the density \eqref{g1:eq:Z-density} and a Laplace saddle point.
\end{conjecture}

The experiment in \cref{g1:fig:Fabius-ratio} is consistent with a slowly
converging logarithmic-square law, but orders below $40$ are far too small to
distinguish the lower-order $\log n\log\log n$ and $\log n$ corrections.

\subsection{Second-order regular variation}

\cref{g1:eq:regular-variation} gives only the first boundary-layer limit.  If
$f(x)=x^\alpha L(x)$ and $L$ belongs to a de Haan second-order class, one should
be able to expand
\[
 \frac{(\Gc_{n,q}f)(A)}{f(Aq^n)}
 =C_q(\alpha)+A_1(Aq^n)C_{q,1}(\alpha,\rho)+o(A_1(Aq^n)),
\]
where $A_1$ is the auxiliary function and the correction constant is a
$q$-Mellin moment of the signed boundary layer $\beta_r^{(q)}$.  This would
provide rigorous error estimators for nonanalytic asymptotic extrapolation.

\begin{problem}[Second-order transfer theorem]
Formulate minimal second-order regular-variation hypotheses and compute the
correction constants explicitly from
$(qz;q)_\infty/(q;q)_\infty$.  Treat the resonant integer cases, where
$C_q(\alpha)=0$, separately; there the correction may become the leading term.
\end{problem}

\subsection{The double-scaling transition $q\uparrow1$, $n\to\infty$}

The fixed-order behavior \eqref{g1:eq:fixed-n-q1} and infinite-order behavior
\eqref{g1:eq:K-modular} do not commute.  A natural transition parameter is
$n(1-q)$ or $nt$ when $q=e^{-t}$.

\begin{problem}[Universal transition kernel]
Determine asymptotics of the row polynomial, its total variation, and the
Peano-kernel density when
\[
 n\to\infty,\qquad t\downarrow0,
 \qquad nt\to\tau\in(0,\infty).
\]
One expects a continuum limit connecting geometric-comb extrapolation to
confluent Hermite/Taylor extrapolation at $A$.  The relevant asymptotics should
involve dilogarithms and a saddle point for finite $q$-products.
\end{problem}

\subsection{Optimal filters under analytic priors}

The Gaussian filter is algebraically unique when exactly $n$ integer modes are
canceled with $n+1$ samples.  With more samples than constraints, stability and
bias can be traded.

\begin{problem}[Minimax geometric extrapolation]
Given $M>n$, choose weights on $A,Aq,\ldots,Aq^M$ that preserve constants and
annihilate $x,\ldots,x^n$, while minimizing either total variation or the worst
error over an analytic coefficient-majorant class.  Compare the optimizer with
truncated, damped, and least-squares modifications of $W_{n,q}$.
\end{problem}

The problem is finite-dimensional and convex for total variation, but its
large-$M,n$ structure is likely governed by weighted extremal polynomials on
the spectral set $\{q^m:m\ge0\}$.

\subsection{Inverse Fabius and nonlinear combs}

The inverse Fabius function has endpoint behavior governed by asymptotic
inversion and Lambert $W$, not a regular Taylor series.  A promising strategy is
to choose a nonlinear sample comb $x_k$ so that the dominant inverse asymptotic
modes become geometric in a transformed coordinate.  Confluent modal filters
from \cref{g1:sec:modal} can then remove the leading Lambert-$W$ transseries terms.

\begin{problem}[Lambert-adapted Newton coordinates]
Find a coordinate $u=u(x)$ in which the inverse Fabius germ admits an expansion
in modes $e^{-\alpha u}u^s$, construct the corresponding geometric Newton
series, and determine whether its interpolation coefficients have a positive
kernel or a controlled boundary layer analogous to \eqref{g1:eq:beta}.
\end{problem}

\subsection{Multivariate geometric combs}

Tensor products immediately give interpolation on
$\{A_1q_1^{k_1}\}\times\cdots\times\{A_dq_d^{k_d}\}$, with products of
Gaussian rows.  More interesting are diagonal or simplex combs with a shared
scale index, which interact with multivariate self-similar distributions and
box splines.  The Dirichlet-average interpretation suggests that multivariate
Peano kernels may again have positive probabilistic representations.

\subsection{Formalization boundary and remaining roadmap}

The live Lean corpus has moved beyond a roadmap for the elementary finite
filter.  The modules \nolinkurl{FiniteQBinomialCore.lean},
\nolinkurl{GeometricLagrangeQBinomial.lean},
\nolinkurl{GeometricLagrangeQMoments.lean}, and
\nolinkurl{GeometricLagrangeCompleteHomogeneous.lean} contain named theorems
for finite $q$-products, Gaussian weights, cancellation/residual formulas, the
complete-homogeneous specialization, and the exact finite row norm.
\nolinkurl{LimitConditionNumber.lean} supplies the fixed-$q$ infinite-product
limit, while the geometric-uniform and Rvachev modules named in the opening
status box cover related probability and Fourier infrastructure.

The report-specific obligations that remain are therefore narrower:
\begin{enumerate}[leftmargin=2em]
\item identify geometric divided differences with iterated Jackson derivatives
under explicit field and differentiability hypotheses;
\item construct the infinite Newton series and prove its local uniform
convergence;
\item formalize confluent spectral zeros and the Puiseux--log remainder bound;
\item prove the reversed-row $\ell^1$ limit and the regular-variation transfer;
\item add Hermite--Genocchi/Dirichlet-average infrastructure, the geometric-knot
B-spline moments, and the perpetuity scaling limit;
\item derive the Fabius boundary expectation and the entire scale-reconstruction
theorem from the already formalized finite core.
\end{enumerate}
Until those items exist as checked declarations, the correspondingly labelled
results here remain manuscript theorems.

\section{Complementary Mellin and superflat theory}
\subsection{A target on the same multiplicative line}

A target $x=aq^\alpha$, where $\alpha$ need not be integral, preserves all
$q$-structure.

\begin{theorem}[Arbitrary geometric target]
\label{g2:thm:general-target}
Choose a branch of $q^\alpha$.  If $q^\alpha$ is not one of
$1,q,\ldots,q^n$, then
\begin{equation}
 \ell_{n,j}(aq^\alpha)
 =(-1)^{n-j}q^{\binom{n-j+1}{2}}
 \frac{(q^{\alpha-n};q)_{n+1}}
 {(1-q^{\alpha-j})(q;q)_j(q;q)_{n-j}}.
 \label{g2:eq:general-target-cardinal}
\end{equation}
At $\alpha=j$ the singularity is removable and the value is $1$; at every
other integral node it is $0$.  For $0<q<1$ and
$\Re\alpha\to+\infty$, the weights tend to the zero-target row
\begin{equation}
 \lambda_{n,j}^{(q)}
 =\ell_{n,j}(0)
 =\frac{(-1)^{n-j}q^{\binom{n-j+1}{2}}}
 {(q;q)_j(q;q)_{n-j}}.
 \label{g2:eq:zero-target-row}
\end{equation}
\end{theorem}

\begin{proof}
Use
$\ell_{n,j}(x)=\Omega_{n+1}(x)/((x-aq^j)\Omega'_{n+1}(aq^j))$,
\eqref{gq:eq:nodal-polynomial} and \eqref{gq:eq:barycentric-weight}.  The product
at the target is
\[
 \prod_{r=0}^n(q^\alpha-q^r)
 =(-1)^{n+1}q^{\binom{n+1}{2}}(q^{\alpha-n};q)_{n+1}.
\]
Extracting $-q^j$ from $q^\alpha-q^j$ leaves exactly
\eqref{g2:eq:general-target-cardinal}.  The limiting assertions follow directly
from the product.
\end{proof}

\begin{remark}[$q$-gamma form]
Since
\[
 (q^{\alpha-n};q)_{n+1}
 =(1-q)^{n+1}\frac{\Gamma_q(\alpha+1)}{\Gamma_q(\alpha-n)},
\]
formula \eqref{g2:eq:general-target-cardinal} is meromorphic in $\alpha$ and can
be rewritten entirely with $q$-gamma functions and the $q$-number
$[\alpha-j]_q$.  This is the natural analytic continuation of the cardinal
row, not merely a formal replacement of factorials by $q$-factorials.
\end{remark}


\section{Mellin symbols and exact residuals}
\label{g2:sec:mellin}

\subsection{Complex powers}

The integer moment law in the source monograph is one specialization of a
stronger identity.  Since dilations act diagonally on every chosen branch of
$x^s$, the full spectral response is available at once.

\begin{theorem}[Exact Mellin symbol]
\label{g2:thm:mellin-symbol}
Let $s\in\ComplexNumbers$, choose a branch of $a^s$ and $q^s$, and set
\begin{equation}
 M_{n,q}(s)=\sum_{k=0}^n\lambda_{n,k}^{(q)}q^{ks}.
 \label{g2:eq:mellin-symbol-definition}
\end{equation}
Then
\begin{align}
 M_{n,q}(s)
 &=\frac{\prod_{j=1}^n(q^s-q^j)}{(q;q)_n}
 \label{g2:eq:mellin-product}\\
 &=q^{ns}\frac{(q^{1-s};q)_n}{(q;q)_n}
 \label{g2:eq:mellin-forward}\\
 &=(-1)^nq^{\binom{n+1}{2}}
   \frac{(q^{s-n};q)_n}{(q;q)_n}.
 \label{g2:eq:mellin-reversed}
\end{align}
Equivalently,
\begin{equation}
 (\Gop_{n,q}x^s)(a)=a^sM_{n,q}(s).
 \label{g2:eq:mellin-eigenvalue}
\end{equation}
The zeros at $s=1,\ldots,n$ are simple modulo the imaginary period
$2\pi\ImaginaryUnit/\log q$.
\end{theorem}

\begin{proof}
Evaluate \eqref{g1:eq:row-polynomial} at $z=q^s$.  Factoring $q^s$ from
each factor gives \eqref{g2:eq:mellin-forward}; factoring $-q^j$ gives
\eqref{g2:eq:mellin-reversed}.
\end{proof}

\begin{corollary}[Natural moments and the first residual]
\label{g2:cor:integer-moments}
For $d\in\NaturalNumbers$,
\begin{equation}
 \sum_{k=0}^n\lambda_{n,k}^{(q)}q^{kd}
 =\begin{cases}
 1,&d=0,\\
 0,&1\le d\le n,\\
 (-1)^nq^{\binom{n+1}{2}}\GaussianBinomial{d-1}{n}q,&d\ge n+1.
 \end{cases}
 \label{g2:eq:integer-moment-law}
\end{equation}
\end{corollary}

\begin{proof}
The first two cases follow from the zeros in
\eqref{g2:eq:mellin-product}.  For $d\ge n+1$,
\[
 (q^{d-n};q)_n=\frac{(q;q)_{d-1}}{(q;q)_{d-n-1}},
\]
and substitution into \eqref{g2:eq:mellin-reversed} gives the Gaussian
coefficient.
\end{proof}

\begin{remark}[Fractional endpoint laws]
For fixed $s$ with $\Re s>0$ and $s\notin\NaturalNumbers$,
\begin{equation}
 M_{n,q}(s)
 \sim q^{ns}\frac{(q^{1-s};q)_\infty}{(q;q)_\infty}
 \qquad(n\to\infty).
 \label{g2:eq:fractional-asymptotic}
\end{equation}
Thus a fractional power $x^s$ is not annihilated, but its extrapolation error
has a completely explicit geometric rate.  Integer powers are exceptional
because the finite product terminates at an exact zero.
\end{remark}

\subsection{Logarithmic powers by spectral differentiation}

Differentiate \eqref{g2:eq:mellin-eigenvalue} with respect to $s$.  This gives
closed responses for singular terms common in endpoint asymptotics.

\begin{proposition}[Powers times logarithms]
\label{g2:prop:log-powers}
For every integer $r\ge0$,
\begin{equation}
 \sum_{k=0}^n\lambda_{n,k}^{(q)}(aq^k)^s
 \bigl(\log a+k\log q\bigr)^r
 =\frac{\partial^r}{\partial s^r}\bigl(a^sM_{n,q}(s)\bigr).
 \label{g2:eq:log-power-response}
\end{equation}
In particular, if $1\le d\le n$, then
\begin{equation}
 \sum_{k=0}^n\lambda_{n,k}^{(q)}(aq^k)^d\log(aq^k)
 =a^dM_{n,q}'(d),
 \label{g2:eq:first-log-response}
\end{equation}
where
\begin{equation}
 M_{n,q}'(d)
 =(-1)^{d-1}(\log q)
 q^{d(2n-d+1)/2}
 \frac{(q;q)_{d-1}(q;q)_{n-d}}{(q;q)_n}.
 \label{g2:eq:mellin-derivative-integer}
\end{equation}
\end{proposition}

\begin{proof}
Termwise differentiation is legitimate because the sum is finite.  At
$s=d$, the factor $q^s-q^d$ is the unique vanishing factor in
\eqref{g2:eq:mellin-product}.  Differentiating that factor and splitting the
remaining product below and above $d$ gives
\begin{align*}
 q^d\log q\prod_{\substack{1\le j\le n\\j\ne d}}(q^d-q^j)
 &=(-1)^{d-1}(\log q)
 q^{d(2n-d+1)/2}(q;q)_{d-1}(q;q)_{n-d}.
\end{align*}
Division by $(q;q)_n$ proves \eqref{g2:eq:mellin-derivative-integer}.
\end{proof}

\begin{remark}
A filter that is exact on polynomials need not suppress logarithmic
corrections attached to those powers.  Formula
\eqref{g2:eq:mellin-derivative-integer} quantifies the leakage exactly.  Higher
$s$-derivatives introduce finite sums of $q$-harmonic terms
$(1-q^j)^{-r}$ and can be organized with complete Bell polynomials.
\end{remark}

\subsection{Mellin-transform interpretation}

For a function on $(0,\infty)$ with Mellin transform
\[
 \MellinTransformOf{f}(z)=\int_0^\infty f(x)x^{z-1}\Differential x,
\]
one has
\[
 \MellinTransformOf{f(q^kx)}(z)=q^{-kz}\MellinTransformOf{f}(z).
\]
Hence, on a common strip of convergence,
\begin{equation}
 \MellinTransformOf{\Gop_{n,q}f}(z)=M_{n,q}(-z)\MellinTransformOf{f}(z).
 \label{g2:eq:mellin-transform-multiplier}
\end{equation}
The geometric Lagrange filter is therefore a finite Mellin multiplier.  Its
spectral zeros are the same integer moments that the interpolation formula
annihilates in physical space.  This is the multiplicative analogue of an
ordinary finite-difference stencil whose Fourier symbol has a prescribed zero
at the origin.


\subsection{The limits \texorpdfstring{$q\to1$}{q to 1} and \texorpdfstring{$q\to0$}{q to 0}}

\begin{proposition}[Taylor limit]
\label{g2:prop:taylor-limit}
For fixed $k$, as $q\to1$,
\[
 D_q^kf(a)\longrightarrow f^{(k)}(a),
 \qquad [k]_q!\longrightarrow k!,
 \qquad \Phi_k(x)\longrightarrow(x-a)^k.
\]
Thus \eqref{gq:eq:q-taylor-interpolant} tends termwise to the ordinary Taylor
polynomial about $a$.
\end{proposition}

\begin{proof}
The first statement follows by repeated use of the standard Jackson derivative
limit on a sufficiently differentiable function; the other two are immediate
from their finite products.
\end{proof}

At the opposite endpoint, $D_0f(x)=(f(x)-f(0))/x$ is a tail-quotient operator.
For $f(x)=\sum c_mx^m$,
\[
 \lim_{q\to0}\frac{D_q^kf(a)}{[k]_q!}
 =\sum_{r\ge0}c_{k+r}a^r,
 \qquad
 \lim_{q\to0}\Phi_k(x)=
 \begin{cases}
 1,&k=0,\\
 (x-a)x^{k-1},&k\ge1.
 \end{cases}
\]
The limiting expansion telescopes.  It is not Hermite interpolation at a
multiple node: the geometric rates retain ordered information about how the
nodes approach zero.  A genuine confluent theory requires a further
renormalization and is listed as \Cref{g2:prob:q0-confluent}.

\subsection{Neville, Richardson, and Romberg recurrences}

\begin{theorem}[Geometric Neville recurrence]
\label{g2:thm:neville}
For $i,m\ge0$, let $P_{i,m}$ be the interpolant characterized by
\[
 P_{i,m}(aq^{i+r})=f(aq^{i+r})\qquad(0\le r\le m).
\]
Then
\begin{equation}
 P_{i,m}(x)=
 \frac{(x-aq^{i+m})P_{i,m-1}(x)
      -(x-aq^i)P_{i+1,m-1}(x)}
      {aq^i-aq^{i+m}}.
 \label{g2:eq:geometric-neville}
\end{equation}
At $x=0$, with $R_i^{(0)}=f(aq^i)$,
\begin{equation}
 R_i^{(m)}
 =\frac{R_{i+1}^{(m-1)}-q^mR_i^{(m-1)}}{1-q^m},
 \qquad
 R_0^{(n)}=(\Gop_{n,q}f)(a).
 \label{g2:eq:richardson-recurrence}
\end{equation}
\end{theorem}

\begin{proof}
The numerator in \eqref{g2:eq:geometric-neville} has degree at most $m$ and takes
the required values at both endpoint nodes and all shared interior nodes;
uniqueness proves the formula.  Substitute $x=0$ and cancel $aq^i$.
\end{proof}

\begin{corollary}[Generalized Richardson cancellation]
\label{g2:cor:richardson}
Suppose approximations at scales $h_i=h_0r^i$ have an expansion
\[
 T(h)=L+c_1h^\rho+c_2h^{2\rho}+\cdots.
\]
Apply \eqref{g2:eq:richardson-recurrence} with $q=r^\rho$ to the data $T(h_i)$.
After $n$ levels, the first $n$ error powers are annihilated exactly.  Romberg
quadrature is the special case in which the trapezoidal error uses even powers
and hence $q=r^2$.
\end{corollary}

\begin{proof}
The recurrence is the degree-$n$ interpolation functional at the point
$0$ on nodes $1,q,\ldots,q^n$.  It fixes the constant mode and annihilates
$q^{ki}$ for $1\le k\le n$, equivalently the error modes
$h^{k\rho}$.  Linearity therefore removes the first $n$ displayed error
powers.  For the trapezoidal expansion only even powers occur, so replacing
$r$ by $r^2$ gives the Romberg case.
\end{proof}

\begin{remark}
This recurrence is often numerically preferable to forming the closed weights.
It also explains the historical term ``deferred approach to the limit'': each
new level eliminates one Mellin eigenmode of the asymptotic error.
\end{remark}


\section{Convergence at an accumulation point and divergence away from it}
\label{g2:sec:superflat}

\subsection{The zero row is a regular summability method}

\begin{theorem}[Regularity at the accumulation point]
\label{g2:thm:regular-summability}
Let $0<q<1$ and let $y_k\to L$.  Then
\begin{equation}
 \sum_{k=0}^n\lambda_{n,k}^{(q)}y_k\longrightarrow L.
 \label{g2:eq:regular-summability}
\end{equation}
Equivalently, if $f$ is continuous at $0$, the interpolating polynomial through
$(aq^k,f(aq^k))$, $0\le k\le n$, satisfies
\begin{equation}
 (\Iop_nf)(0)\longrightarrow f(0).
 \label{g2:eq:continuous-zero-convergence}
\end{equation}
\end{theorem}

\begin{proof}
The rows reproduce constants by \eqref{g2:eq:integer-moment-law}; their
$\ell^1$ norms are bounded by $K(q)$ by \Cref{g1:thm:row-norm}; and for each fixed
$k$,
\[
 |\lambda_{n,k}^{(q)}|
 =\frac{q^{\binom{n-k+1}{2}}}
 {(q;q)_k(q;q)_{n-k}}\longrightarrow0.
\]
These are precisely the Toeplitz regularity conditions.  For completeness,
subtract $L$, split the sum at a fixed $K$, make the tail small using the
bounded row norm, and then let $n\to\infty$ in the finite head.
\end{proof}

\subsection{The limiting profile viewed from the finest tooth}

Reindex the row by $r=n-k$.  For each fixed $r$,
\begin{equation}
 \lambda_{n,n-r}^{(q)}
 \longrightarrow
 \mu_r^{(q)}
 :=\frac{(-1)^rq^{\binom{r+1}{2}}}
 {(q;q)_\infty(q;q)_r}.
 \label{g2:eq:limiting-row-profile}
\end{equation}
Euler's infinite theorem gives the generating function
\begin{equation}
 \sum_{r=0}^\infty\mu_r^{(q)}z^r
 =\frac{(qz;q)_\infty}{(q;q)_\infty},
 \qquad
 \sum_{r=0}^\infty\mu_r^{(q)}=1.
 \label{g2:eq:limiting-profile-generating}
\end{equation}
Thus the mass does not remain attached to any fixed physical node; it travels
with the finest scale $aq^n$ and approaches a signed, absolutely summable
profile in the reversed index.  This is the mechanism behind
\Cref{g2:thm:regular-summability}.

\subsection{Power, Puiseux, and logarithmic endpoint expansions}

Suppose a germ has a convergent or asymptotic expansion
\[
 f(x)=L+\sum_{r}c_rx^{\alpha_r}(\log x)^{m_r},
 \qquad \Re\alpha_r>0.
\]
Theorems \ref{g2:thm:mellin-symbol} and \ref{g2:prop:log-powers} give the filter of
each term exactly.  In particular, a nonintegral leading exponent $\alpha$
produces the rate $q^{n\alpha}$ times an explicit $q$-product, while an
integral polynomial jet is deleted once $n$ reaches its degree.  The filter is
therefore an endpoint asymptotic microscope: its decay rate distinguishes
fractional regularity from an analytic integer jet.

\begin{proposition}[A simple Hölder estimate]
\label{g2:prop:holder-filter}
If $|f(x)-f(0)|\le C|x|^\alpha$ on the comb, with $\alpha>0$, then
\begin{equation}
 |(\Gop_{n,q}f)(a)-f(0)|
 \le C|a|^\alpha
 \sum_{k=0}^n|\lambda_{n,k}^{(q)}|q^{k\alpha}.
 \label{g2:eq:holder-bound}
\end{equation}
For fixed $q$ and $\alpha$, the right side is $O(q^{n\alpha})$.
\end{proposition}

\begin{proof}
The first statement is the triangle inequality.  Reindex $r=n-k$.  Since
$(q;q)_m\ge(q;q)_\infty$ for every $m$, one has
\begin{align*}
 q^{-n\alpha}
 \sum_{k=0}^n|\lambda_{n,k}^{(q)}|q^{k\alpha}
 &=\sum_{r=0}^n
 \frac{q^{\binom{r+1}{2}-r\alpha}}
 {(q;q)_{n-r}(q;q)_r}\\
 &\le \frac1{(q;q)_\infty^2}
 \sum_{r=0}^\infty q^{\binom{r+1}{2}-r\alpha}<\infty.
\end{align*}
The last series converges for every fixed $\alpha$ because its exponent is
quadratic in $r$.  This uniform bound proves the stated $O(q^{n\alpha})$
estimate.  Dominated convergence also gives the sharper limit
\[
 q^{-n\alpha}
 \sum_{k=0}^n|\lambda_{n,k}^{(q)}|q^{k\alpha}
 \longrightarrow
 \sum_{r=0}^\infty|\mu_r^{(q)}|q^{-r\alpha}.
\]
\end{proof}

\subsection{Superflat data force divergence off the comb}

The favorable result at $0$ has a nearly opposite counterpart everywhere
else.  It is useful to formulate it directly for data, without assuming that
they initially came from a function.

\begin{definition}
A sequence $y_j\to L$ is \emph{$q$-superflat at its limit} if
\begin{equation}
 y_j-L=o(q^{mj})\qquad(j\to\infty)
 \label{g2:eq:q-superflat}
\end{equation}
for every $m\in\NaturalNumbers$.
\end{definition}

\begin{theorem}[Exact convergence set for nontrivial superflat data]
\label{g2:thm:superflat-divergence}
Let $0<q<1$, $a\ne0$, and let $y_j\to L$ be $q$-superflat.  Assume
$y_j\ne L$ for infinitely many $j$.  Let
\begin{equation}
 P_n(x)=\sum_{k=0}^n A_k\Phi_k(x)
 \label{g2:eq:superflat-newton}
\end{equation}
be the polynomial interpolating $P_n(aq^j)=y_j$ for $0\le j\le n$.  Then
\begin{equation}
 \limsup_{k\to\infty}|A_k|^{1/k}=\infty.
 \label{g2:eq:superexponential-A}
\end{equation}
Moreover:
\begin{enumerate}
\item $P_n(0)\to L$;
\item for each $m\ge0$, $P_n(aq^m)=y_m$ once $n\ge m$;
\item for every $x\in\ComplexNumbers\setminus(\{0\}\cup\{aq^m:m\ge0\})$, the increments
      $A_n\Phi_n(x)$ fail to tend to zero, and the sequence $P_n(x)$ diverges;
      indeed it is unbounded along a subsequence.
\end{enumerate}
Thus the convergence set is exactly the geometric comb together with its
accumulation point.
\end{theorem}

\begin{proof}
Subtract $L$; this changes only $A_0$ and reduces to $L=0$.  Suppose, contrary
to \eqref{g2:eq:superexponential-A}, that $|A_k|\le CR^{-k}$ for some $R>0$.
For every $0<r<R$, the estimate
\[
 \sup_{|z|\le r}|\Phi_k(z)|
 \le r^k\prod_{j=0}^\infty(1+|a|q^j/r)
\]
shows that $g(z)=\sum_{k\ge0}A_k\Phi_k(z)$ is analytic on $|z|<R$.
For every sufficiently large $j$, $|aq^j|<R$, and all terms with $k>j$
vanish at $aq^j$.  Hence
\[
 g(aq^j)=P_j(aq^j)=y_j.
\]
If the first nonzero Taylor coefficient of $g$ at $0$ had order $m$, then
$g(aq^j)/(aq^j)^m$ would tend to a nonzero constant.  This contradicts
$q$-superflatness.  Therefore every Taylor coefficient vanishes, so $g$ is
identically zero near $0$, contradicting the assumption that infinitely many
$y_j$ are nonzero.  This proves \eqref{g2:eq:superexponential-A}.

Assertion 1 is \Cref{g2:thm:regular-summability}, and assertion 2 is interpolation.
For a point $x$ in assertion 3,
\begin{equation}
 \Phi_n(x)=x^n\prod_{r=0}^{n-1}\left(1-\frac{aq^r}{x}\right).
 \label{g2:eq:basis-fixed-x}
\end{equation}
The infinite product tends to $(a/x;q)_\infty\ne0$, because $x$ is not a
node.  Hence $|\Phi_n(x)|^{1/n}\to|x|>0$.  Together with
\eqref{g2:eq:superexponential-A}, this implies
\(
 \limsup|A_n\Phi_n(x)|^{1/n}=\infty
\); the increments do not tend to zero.  Finally,
$|P_n-P_{n-1}|=|A_n\Phi_n|$ is unbounded along a subsequence, so at least one
of each adjacent pair $P_n,P_{n-1}$ is unbounded along a subsequence.
\end{proof}

\begin{remark}[A precise geometric Runge phenomenon]
This is not the classical Runge phenomenon caused by equispaced nodes on a
fixed interval.  The comb is not dense away from $0$.  Superflat nonzero data
cannot be the trace of a nonzero analytic germ at the accumulation point;
the interpolation coefficients must therefore leave every exponential class.
The instability is an analytic-continuation obstruction encoded in a
Vandermonde problem.
\end{remark}

\begin{remark}[Derivative-enhanced interpolation]
Adding finitely many Hermite conditions at any fixed collection of nonzero
teeth does not remove the obstruction: an exponentially bounded coefficient
sequence would still define an analytic germ with a superflat nonzero trace.
A successful regularization must either change the approximation space, use a
number of derivative conditions growing with $n$, or explicitly build
flatness into the basis.  Rvachev atoms are natural candidates for the last
option.
\end{remark}


\section{The geometric-uniform family}
\label{g2:sec:geometric-uniform}

\subsection{Probability law, support, and refinement}

Let $U_0,U_1,\ldots$ be independent uniform random variables on $[0,1]$ and
set
\begin{equation}
 X_q=(1-q)\sum_{j=0}^\infty q^jU_j,
 \qquad 0<q<1.
 \label{g2:eq:Xq-definition}
\end{equation}
This is a geometric member of the infinite-convolution class studied by
Jessen and Wintner~\cite{JessenWintner1935}.  The coefficients are positive
and sum to $1$, so $0\le X_q\le1$.  Reflection
of every $U_j$ shows
\begin{equation}
 X_q\stackrel d=1-X_q.
 \label{g2:eq:Xq-symmetry}
\end{equation}
Let $F_q(x)=\Probability(X_q\le x)$, extended by $0$ to $x\le0$ and by $1$ to
$x\ge1$.  Removing the first coordinate gives the affine fixed-point law
\begin{equation}
 X_q\stackrel d=(1-q)U+qX_q',
 \label{g2:eq:Xq-fixed-point}
\end{equation}
where $X_q'$ is an independent copy.  Consequently
\begin{align}
 F_q(x)
 &=\int_0^1F_q\!\left(\frac{x-(1-q)u}{q}\right)\Differential u
 \label{g2:eq:Fq-integral-equation}\\
 &=\frac q{1-q}
 \int_{(x-(1-q))/q}^{x/q}F_q(t)\Differential t.
 \label{g2:eq:Fq-integral-equation-2}
\end{align}
Differentiation gives
\begin{equation}
 F_q'(x)=\frac{F_q(x/q)-F_q((x-(1-q))/q)}{1-q}.
 \label{g2:eq:Fq-derivative}
\end{equation}

\begin{proposition}[Smooth compactly supported density]
\label{g2:prop:uq-smooth}
The centered random variable
\begin{equation}
 Z_q=2X_q-1=(1-q)\sum_{j=0}^\infty q^jV_j,
 \qquad V_j\sim\operatorname{Unif}[-1,1],
 \label{g2:eq:Zq-definition}
\end{equation}
has an even $C^\infty$ density $u_q$ supported on $[-1,1]$.  With the Fourier
convention
$\FourierTwoPi{f}(\xi)=\int_\RealNumbers f(x)\EulerE^{-2\pi\ImaginaryUnit\xi x}\Differential x$,
\begin{equation}
 \FourierTwoPi{u_q}(\xi)=\prod_{j=0}^\infty
 \SincPi\!\left(2(1-q)q^j\xi\right),
 \qquad
 \SincPi(t)=\frac{\sin(\pi t)}{\pi t}.
 \label{g2:eq:uq-fourier-product}
\end{equation}
Moreover
\begin{equation}
 u_q(z)=\frac12F_q'\!\left(\frac{z+1}{2}\right).
 \label{g2:eq:uq-Fq-relation}
\end{equation}
\end{proposition}

\begin{proof}
The characteristic function factors by independence and gives
\eqref{g2:eq:uq-fourier-product}.  To see rapid decay, use
$|\SincPi(t)|\le\min(1,1/(\pi|t|))$ on the first
$N\asymp\log|\xi|/|\log q|$ factors.  The logarithm of the resulting product
is $-c_q(\log|\xi|)^2+O_q(\log|\xi|)$, hence the product decreases faster than
any inverse power.  Fourier inversion then gives a $C^\infty$ density.  The
support and evenness follow from \eqref{g2:eq:Zq-definition}; the density change
of variables gives \eqref{g2:eq:uq-Fq-relation}.
\end{proof}

\begin{remark}
The repository's formal development treats a still broader geometric-uniform
law, including negative ratios and uniqueness of the affine distributional
fixed point.  The positive-ratio case above is the one naturally matched to a
one-sided interpolation comb.
\end{remark}

\subsection{Moments and Bernoulli cumulants}

Write
\begin{equation}
 \mu_n(q)=\Expectation X_q^n.
 \label{g2:eq:muq-definition}
\end{equation}

\begin{theorem}[Moment recurrence]
\label{g2:thm:moment-recurrence}
The moments satisfy $\mu_0(q)=1$ and, for $n\ge1$,
\begin{equation}
 (1-q^n)\mu_n(q)
 =\sum_{k=0}^{n-1}\binom nk
 q^k(1-q)^{n-k}\frac{\mu_k(q)}{n-k+1}.
 \label{g2:eq:general-moment-recurrence}
\end{equation}
At $q=1/2$ this reduces to
\begin{equation}
 (2^n-1)\mu_n(1/2)
 =\sum_{k=0}^{n-1}\binom nk\frac{\mu_k(1/2)}{n-k+1}.
 \label{g2:eq:dyadic-moment-recurrence}
\end{equation}
\end{theorem}

\begin{proof}
Raise \eqref{g2:eq:Xq-fixed-point} to the $n$th power, expand, and use
$\Expectation U^r=1/(r+1)$.  Move the $k=n$ term to the left.  At $q=1/2$ every factor
$q^k(1-q)^{n-k}$ equals $2^{-n}$, and multiplication by $2^n$ gives
\eqref{g2:eq:dyadic-moment-recurrence}.
\end{proof}

The centered moment generating function is
\begin{equation}
 \Mq(z)=\Expectation\EulerE^{zZ_q}
 =\prod_{j=0}^\infty
 \frac{\sinh((1-q)q^jz)}{(1-q)q^jz}.
 \label{g2:eq:uq-mgf}
\end{equation}
Using
\[
 \log\frac{\sinh z}{z}
 =\sum_{m\ge1}\frac{2^{2m-1}B_{2m}}{m(2m)!}z^{2m},
\]
we obtain the even cumulants
\begin{equation}
 \kappa_{2m}(Z_q)
 =\frac{2^{2m}B_{2m}}{2m}
   \frac{(1-q)^{2m}}{1-q^{2m}},
 \qquad
 \kappa_{2m+1}(Z_q)=0.
 \label{g2:eq:uq-cumulants}
\end{equation}
Thus the same geometric series that produces the interpolation denominators
also produces the cumulant denominators.  At $q=1/2$,
\begin{equation}
 \kappa_{2m}(Z_{1/2})
 =\frac{B_{2m}}{2m(1-4^{-m})}.
 \label{g2:eq:up-cumulants}
\end{equation}
Moments can be recovered from these cumulants by complete Bell polynomials,
providing a second exact route parallel to \eqref{g2:eq:dyadic-moment-recurrence}.

\subsection{The left boundary layer}

For $0\le x\le1-q$, the second term in \eqref{g2:eq:Fq-derivative} vanishes, so
\begin{equation}
 F_q'(x)=\frac1{1-q}F_q(x/q).
 \label{g2:eq:left-boundary-first}
\end{equation}
Iteration gives the geometric derivative tower.

\begin{lemma}[Boundary derivative tower]
\label{g2:lem:boundary-tower}
If $m\ge1$ and $0\le x\le(1-q)q^{m-1}$, then
\begin{equation}
 F_q^{(m)}(x)
 =(1-q)^{-m}q^{-\binom m2}F_q(x/q^m).
 \label{g2:eq:boundary-derivative-tower}
\end{equation}
\end{lemma}

\begin{proof}
Induct on $m$.  The case $m=1$ is
\eqref{g2:eq:left-boundary-first}.  Differentiating the induction hypothesis
contributes a factor $q^{-1}$, and the domain assumption ensures that the next
use of \eqref{g2:eq:left-boundary-first} remains in the left boundary layer.
\end{proof}

\subsection{Comb samples are normalized moments}

The ratio $q=1/2$ is geometrically critical: the first uniform block has
length $1-q$, while the tail has length $q$, and these are equal exactly at
$q=1/2$.  For every $q\le1/2$, however, the tooth $q^n$ still lies inside the
$n$th iterated boundary layer, which is enough for the following formula.

\begin{theorem}[Geometric comb--moment identity]
\label{g2:thm:comb-moment}
For $0<q\le1/2$ and every $n\ge0$,
\begin{equation}
 F_q(q^n)
 =\frac{q^{\binom{n+1}{2}}}{(1-q)^n n!}\,\mu_n(q).
 \label{g2:eq:comb-moment}
\end{equation}
\end{theorem}

\begin{proof}
The case $n=0$ is $F_q(1)=1$.  For $n\ge1$, compact support and smoothness
imply that $F_q$ and all its derivatives vanish at $0$.  Since
$q^n\le(1-q)q^{n-1}$ when $q\le1/2$, Taylor's integral formula and
\Cref{g2:lem:boundary-tower} give
\begin{align*}
 F_q(q^n)
 &=\frac1{(n-1)!}\int_0^{q^n}(q^n-t)^{n-1}F_q^{(n)}(t)\Differential t\\
 &=\frac{(1-q)^{-n}q^{-\binom n2}}{(n-1)!}
   \int_0^{q^n}(q^n-t)^{n-1}F_q(t/q^n)\Differential t\\
 &=\frac{q^{n^2-\binom n2}}{(1-q)^n(n-1)!}
   \int_0^1(1-u)^{n-1}F_q(u)\Differential u.
\end{align*}
By Fubini,
\[
 \int_0^1(1-u)^{n-1}F_q(u)\Differential u
 =\frac1n\Expectation(1-X_q)^n.
\]
Symmetry \eqref{g2:eq:Xq-symmetry} changes the last moment to $\mu_n(q)$, and
$n^2-\binom n2=\binom{n+1}{2}$.
\end{proof}

\begin{remark}[What fails for $q>1/2$]
The formula is not expected to survive unchanged.  The point $q^n$ lies beyond
the pure left boundary layer, so the reflected term in
\eqref{g2:eq:Fq-derivative} enters before the $n$th differentiation.  The
resulting correction should be indexed by affine words recording boundary
crossings.  Deriving a closed word-sum or automaton formula is
\Cref{g2:prob:qgt-half}.
\end{remark}

\subsection{The classical Fabius and Rvachev normalizations}

Put
\[
 F=F_{1/2},\qquad X=X_{1/2},\qquad u=u_{1/2}.
\]
Then $F$ is the bounded Fabius function.  On $0\le x\le1/2$,
\[
 F'(x)=2F(2x),
 \qquad
 F(1-x)=1-F(x).
\]
The centered density $u$ is Rvachev's up-function in the probability
normalization.  From \eqref{g2:eq:uq-Fq-relation},
\begin{equation}
 u(z)=
 \begin{cases}
 F(z+1),&-1\le z\le0,\\
 F(1-z),&0\le z\le1,\\
 0,&|z|\ge1.
 \end{cases}
 \label{g2:eq:up-Fabius-piecewise}
\end{equation}
Its Fourier transform is the canonical infinite sinc product
\begin{equation}
 \FourierTwoPi{u}(\xi)=\prod_{j=0}^\infty\SincPi(\xi/2^j).
 \label{g2:eq:up-fourier-classical}
\end{equation}
The construction and these identities agree with the classical treatments of
Fabius~\cite{Fabius1966} and Arias de Reyna~\cite{Arias2017Definition}, and
with the normalization used in the repository's primary exposition
\cite{ProveItFabius}.

Write $\mu_n:=\mu_n(1/2)$ and introduce the Mersenne product
\begin{equation}
 \mathfrak M_n:=\prod_{r=1}^n(2^r-1),
 \qquad \mathfrak M_0:=1.
 \label{g2:eq:mersenne-definition}
\end{equation}

\begin{proposition}[Mersenne form of the dyadic origin row]
\label{g2:prop:dyadic-mersenne-row}
For every $n\ge0$,
\begin{align}
 (1/2;1/2)_n
 &=2^{-\binom{n+1}{2}}\mathfrak M_n,
 \label{g2:eq:qpoch-mersenne}\\
 \lambda_{n,k}^{(1/2)}
 &=(-1)^{n-k}{2^{\binom{k+1}{2}}
 \over\mathfrak M_k\mathfrak M_{n-k}}
 \qquad(0\le k\le n).
 \label{g2:eq:dyadic-lambda-mersenne}
\end{align}
\end{proposition}

\begin{proof}
In the Pochhammer product, write
$1-2^{-r}=2^{-r}(2^r-1)$ and multiply the powers $2^{-r}$ for
$1\le r\le n$.  Substitution of that identity at $k$ and $n-k$ into
\eqref{g1:eq:lambda} cancels the remaining powers of two and yields the
second formula.
\end{proof}


\subsection{The endpoint extrapolant as a Mersenne--moment convolution}

Let $P_n$ interpolate $F$ at
$1,1/2,\ldots,2^{-n}$ and put $E_n=P_n(0)$.

\begin{theorem}[Exact endpoint convolution]
\label{g2:thm:fabius-endpoint-convolution}
For every $n\ge0$,
\begin{equation}
 E_n
 =\sum_{k=0}^n(-1)^{n-k}
 \frac{2^k\mu_k}
 {k!\,\mathfrak M_k\mathfrak M_{n-k}}.
 \label{g2:eq:fabius-endpoint-convolution}
\end{equation}
The first values are
\[
 1,\quad 0,\quad -\frac{13}{27},\quad \frac19,\quad
 \frac{4418}{637875},\quad -\frac{3218}{1318275},\ldots.
\]
Despite the initially nonmonotone behavior, $E_n\to0$.
\end{theorem}

\begin{proof}
Substitute \eqref{gq:eq:fabius-inverse-dyadic} and
\eqref{g2:eq:dyadic-lambda-mersenne} into
$E_n=\sum_k\lambda_{n,k}^{(1/2)}F(2^{-k})$.  The powers combine as
$2^{\binom{k+1}{2}-\binom k2}=2^k$.  Convergence follows independently from
\Cref{g2:thm:regular-summability}.
\end{proof}

The same mechanism yields a quantitative bound throughout the
geometric-uniform family.

\begin{theorem}[Quadratic endpoint bound]
\label{g2:thm:quadratic-endpoint-bound}
Let $0<q\le1/2$, and let $P_{n,q}$ interpolate $F_q$ at
$1,q,\ldots,q^n$.  Then
\begin{equation}
 |P_{n,q}(0)|
 \le
 \frac{\exp((1-q)^{-1})}{(q;q)_\infty^2}
 q^{\binom{n+1}{2}-\Floor{n^2/4}}.
 \label{g2:eq:quadratic-endpoint-bound}
\end{equation}
In particular the convergence is at least as fast as
$q^{n^2/4+O(n)}$, and hence faster than every fixed geometric rate.
\end{theorem}

\begin{proof}
Use the exact sample formula \eqref{g2:eq:comb-moment}, the bound
$0\le\mu_k(q)\le1$, and \eqref{g2:eq:zero-target-row}:
\begin{align*}
 |P_{n,q}(0)|
 &\le\sum_{k=0}^n
 \frac{q^{\binom{n-k+1}{2}+\binom{k+1}{2}}}
 {(q;q)_k(q;q)_{n-k}}
 \frac{(1-q)^{-k}}{k!}.
\end{align*}
The exponent is
\[
 \binom{n+1}{2}-k(n-k)
 \ge\binom{n+1}{2}-\Floor{n^2/4}.
\]
Also $(q;q)_k,(q;q)_{n-k}\ge(q;q)_\infty$.  Extract the smallest power of
$q$ and majorize the remaining sum by
$\sum_{k\ge0}(1-q)^{-k}/k!=\exp((1-q)^{-1})$.
\end{proof}

\subsection{Newton coefficients as a Gaussian moment transform}

Write the Newton form
\begin{equation}
 P_n(x)=\sum_{k=0}^n A_k\prod_{r=0}^{k-1}(x-2^{-r}).
 \label{g2:eq:fabius-newton-form}
\end{equation}
The coefficient $A_k$ is independent of the final interpolation degree once
the first $k+1$ nodes are present.

\begin{theorem}[General geometric-uniform moment transform]
\label{g2:thm:general-moment-transform}
For $0<q\le1/2$, let $A_n(q)$ be the Newton coefficient for the data
$F_q(1),F_q(q),\ldots$.  Then
\begin{equation}
 A_n(q)=\frac1{(q;q)_n}
 \sum_{j=0}^n(-1)^j
 \GaussianBinomial nj{q^{-1}}
 \left(\frac q{1-q}\right)^j
 \frac{\mu_j(q)}{j!}.
 \label{g2:eq:general-Gaussian-moment-transform}
\end{equation}
The transform is triangular and invertible.  If
$B_k=(q;q)_kA_k(q)$ and $Q=q^{-1}$, then
\begin{equation}
 \frac{\mu_n(q)}{n!}
 =\left(\frac{1-q}{q}\right)^n
 \sum_{k=0}^n(-1)^kQ^{\binom{n-k}{2}}
 \GaussianBinomial nkQ B_k.
 \label{g2:eq:inverse-Gaussian-moment-transform}
\end{equation}
\end{theorem}

\begin{proof}
Insert \eqref{g2:eq:comb-moment} into the explicit divided difference
\eqref{gq:eq:geometric-divided-difference} with $c=1$.  The powers simplify by
\[
 q^{-jn+\binom{j+1}{2}+\binom{j+1}{2}}(1-q)^{-j}
 =q^{-j(n-j)}\left(\frac q{1-q}\right)^j,
\]
and reciprocity
$q^{-j(n-j)}\GaussianBinomial njq=\GaussianBinomial nj{q^{-1}}$ proves
\eqref{g2:eq:general-Gaussian-moment-transform}.
For inversion, factor the transform into a diagonal matrix and the Gaussian
Pascal matrix, and use the mutually inverse basis transforms in
\Cref{gq:thm:gaussian-pascal-basis} with base $Q$; direct multiplication
gives \eqref{g2:eq:inverse-Gaussian-moment-transform}.
\end{proof}

\begin{corollary}[Dyadic Gaussian and Mersenne forms]
\label{g2:cor:dyadic-moment-transform}
For the Fabius function,
\begin{align}
 A_n
 &=\frac1{(1/2;1/2)_n}
 \sum_{j=0}^n(-1)^j\GaussianBinomial nj2\frac{\mu_j}{j!}
 \label{g2:eq:fabius-Gaussian-transform}\\
 &=2^{\binom{n+1}{2}}
 \sum_{j=0}^n
 \frac{(-1)^j\mu_j}
 {j!\,\mathfrak M_j\mathfrak M_{n-j}}.
 \label{g2:eq:fabius-Mersenne-transform}
\end{align}
Moreover
\begin{equation}
 E_n-E_{n-1}=(-1)^n2^{-\binom n2}A_n
 \qquad(n\ge1).
 \label{g2:eq:endpoint-increment}
\end{equation}
\end{corollary}

\begin{proof}
Set $q=1/2$ in \eqref{g2:eq:general-Gaussian-moment-transform}.  Then
$q^{-1}=2$ and $q/(1-q)=1$.  Use
\(
 \GaussianBinomial nj2=\mathfrak M_n/(\mathfrak M_j\mathfrak M_{n-j})
\)
and \eqref{g2:eq:qpoch-mersenne} for the second form.  Finally, the difference
between two successive Newton partial sums at zero is
$A_n\Phi_n(0)=(-1)^n2^{-\binom n2}A_n$.
\end{proof}

\subsection{Generating-function factorization}

Define
\begin{equation}
 B(z)=\sum_{r\ge0}\frac{z^r}{\mathfrak M_r}
 =(-z/2;1/2)_\infty,
 \qquad
 H(z)=\sum_{r\ge0}\frac{(-1)^r\mu_r}{r!\,\mathfrak M_r}z^r.
 \label{g2:eq:BH-generating}
\end{equation}
The equality for $B$ is Euler's infinite $q$-binomial theorem.  The two exact
convolutions above are equivalent to
\begin{align}
 \sum_{n\ge0}2^{-\binom{n+1}{2}}A_nz^n
 &=B(z)H(z),
 \label{g2:eq:newton-generating-factorization}\\
 \sum_{n\ge0}(-1)^nE_nz^n
 &=B(z)H(2z).
 \label{g2:eq:endpoint-generating-factorization}
\end{align}
These identities exhibit the Fabius Newton coefficients as a $q$-Borel-type
transform of the moment generating function multiplied by a universal Euler
product.  They may be useful for saddle-point analysis of $A_n$.

\subsection{Exact divergence of the Fabius interpolants}

The Fabius function is positive on $(0,1]$ and flat at $0$:
$F(x)=o(x^m)$ for every $m$.  Therefore its dyadic samples satisfy the
hypotheses of \Cref{g2:thm:superflat-divergence}.

\begin{corollary}[Exact convergence set for Fabius comb interpolation]
\label{g2:cor:fabius-convergence-set}
Let $P_n$ interpolate $F$ at $1,1/2,\ldots,2^{-n}$.  Then
\begin{equation}
 P_n(x)\text{ converges as }n\to\infty
 \quad\Longleftrightarrow\quad
 x\in\{0,1,1/2,1/4,\ldots\}.
 \label{g2:eq:fabius-convergence-set}
\end{equation}
At $0$ the limit is $0$ and obeys the quadratic bound
\eqref{g2:eq:quadratic-endpoint-bound}; at a sampled point $2^{-m}$ the sequence
is eventually equal to $F(2^{-m})$; at every other nonzero complex point it is
unbounded along a subsequence.
\end{corollary}

\begin{proof}
Apply \Cref{g2:thm:superflat-divergence} with $a=1$, $q=1/2$, and
$y_j=F(2^{-j})$.
\end{proof}

\begin{corollary}[A rigorous upper scale for coefficient growth]
\label{g2:cor:coefficient-upper}
The Fabius Newton coefficients satisfy
\begin{equation}
 |A_n|\le2^{n^2/4+O(n)}.
 \label{g2:eq:coefficient-upper}
\end{equation}
On the other hand, $\limsup|A_n|^{1/n}=\infty$.
\end{corollary}

\begin{proof}
Combine \eqref{g2:eq:endpoint-increment} with the bounds for $E_n$ and $E_{n-1}$
in \Cref{g2:thm:quadratic-endpoint-bound}.  The lower qualitative statement is
part of \Cref{g2:thm:superflat-divergence}.
\end{proof}

\begin{conjecture}[Quadratic coefficient exponent]
\label{g2:conj:quarter-exponent}
For the Fabius coefficients in \eqref{g2:eq:fabius-newton-form},
\begin{equation}
 \lim_{n\to\infty}\frac{\log_2|A_n|}{n^2}=\frac14.
 \label{g2:eq:quarter-exponent}
\end{equation}
More generally, for $0<q\le1/2$,
\begin{equation}
 \lim_{n\to\infty}
 \frac{\log|A_n(q)|}{n^2\log(1/q)}=\frac14.
 \label{g2:eq:quarter-exponent-general}
\end{equation}
A weaker and probably more robust version replaces the limit by a limsup.  The
Gaussian coefficient in \eqref{g2:eq:general-Gaussian-moment-transform} is largest
near $j=n/2$, where $q^{-j(n-j)}=q^{-n^2/4+O(1)}$; the open issue is to control
alternating cancellation and the subquadratic moment factor.  The order-$80$
companion computation recorded in
\path{assets/companion-evidence/geometric_comb_interpolation_report-3/NUMERICAL_SUMMARY.txt}
is consistent with
\eqref{g2:eq:quarter-exponent}.
\end{conjecture}


\section{The Rvachev Fourier product through the same filter}
\label{g2:sec:rvachev-filter}

\subsection{Tail products along a geometric comb}

Put
\begin{equation}
 s_q(\xi)=\SincPi(2(1-q)\xi),
 \qquad
 S_{q,k}(\xi)=\prod_{j=0}^{k-1}s_q(q^j\xi),
 \quad S_{q,0}=1.
 \label{g2:eq:prefix-products}
\end{equation}
The infinite product \eqref{g2:eq:uq-fourier-product} satisfies
\begin{equation}
 \FourierTwoPi{u_q}(q^k\xi)
 =\frac{\FourierTwoPi{u_q}(\xi)}{S_{q,k}(\xi)}.
 \label{g2:eq:tail-product-ratio}
\end{equation}
Thus values of $\FourierTwoPi{u_q}$ on a geometric frequency comb are exactly the
tails of its defining product.

\begin{theorem}[Reciprocal sinc-product filter]
\label{g2:thm:reciprocal-product}
Let $0<q<1$ and suppose $\FourierTwoPi{u_q}(\xi)\ne0$.  Define
\begin{equation}
 R_{n,q}(\xi)=\sum_{k=0}^n
 \frac{\lambda_{n,k}^{(q)}}{S_{q,k}(\xi)}.
 \label{g2:eq:reciprocal-approximant}
\end{equation}
Then
\begin{equation}
 R_{n,q}(\xi)\longrightarrow\frac1{\FourierTwoPi{u_q}(\xi)}.
 \label{g2:eq:reciprocal-limit}
\end{equation}
More precisely,
\begin{equation}
 \AbsoluteValue{R_{n,q}(\xi)-\frac1{\FourierTwoPi{u_q}(\xi)}}
 \le
 \frac{q^{\binom{n+1}{2}}}{|\FourierTwoPi{u_q}(\xi)|(q;q)_n}
 \EulerE^{2\pi|\xi|}\frac{(2\pi|\xi|)^{n+1}}{(n+1)!}.
 \label{g2:eq:reciprocal-error-bound}
\end{equation}
\end{theorem}

\begin{proof}
Apply the geometric filter to the entire function of $t$
\[
 h_\xi(t)=\FourierTwoPi{u_q}(t\xi)=\Expectation\EulerE^{-2\pi\ImaginaryUnit t\xi Z_q}.
\]
Since $h_\xi(0)=1$, \Cref{g2:thm:regular-summability} already gives
$\sum_k\lambda_{n,k}h_\xi(q^k)\to1$.  By
\eqref{g2:eq:tail-product-ratio}, this sum is
$\FourierTwoPi{u_q}(\xi)R_{n,q}(\xi)$, proving the limit.
For the quantitative bound, the Taylor coefficients of $h_\xi$ satisfy
$|c_m|\le(2\pi|\xi|)^m/m!$ because $|Z_q|\le1$.  Formula
\eqref{g1:eq:exact-analytic-tail} with $A=1$ and
$\GaussianBinomial{n+r}{n}q\le1/(q;q)_n$ gives
\[
 |\FourierTwoPi{u_q}(\xi)R_{n,q}(\xi)-1|
 \le\frac{q^{\binom{n+1}{2}}}{(q;q)_n}
 \sum_{m=n+1}^\infty\frac{(2\pi|\xi|)^m}{m!}.
\]
The exponential-tail bound
$\sum_{m\ge n+1}t^m/m!\le\EulerE^t t^{n+1}/(n+1)!$ completes the proof.
\end{proof}

\begin{remark}[Zeros and regularized reciprocals]
The hypothesis excludes the union of zero lattices
$2(1-q)q^j\xi\in\IntegerNumbers\setminus\{0\}$.  Near a zero, the individual prefix
reciprocals are badly conditioned.  Multiplying by the known vanishing factor,
or differentiating the filtered identity to reconstruct the finite part,
should produce stable formulas for residues and zero multiplicities; see
\Cref{g2:prob:zero-regularization}.
\end{remark}

\subsection{Moment interpretation of the residual}

Expanding the characteristic function gives
\begin{equation}
 \FourierTwoPi{u_q}(t\xi)
 =\sum_{m=0}^\infty
 \frac{(-2\pi\ImaginaryUnit\xi)^m}{m!}\Expectation Z_q^m\,t^m.
 \label{g2:eq:uq-characteristic-series}
\end{equation}
Substitution into \eqref{g1:eq:exact-analytic-tail} expresses the error in
\eqref{g2:eq:reciprocal-limit} as a Gaussian transform of the centered moments.
Odd moments vanish, so the actual leading term is the first even degree above
$n$.  The reciprocal product identity is therefore the Fourier counterpart of
the Fabius moment transform: the same filter reads moments from the CDF comb
and cancels moments in the characteristic-function comb.


\subsection{The additive Thue--Morse companion}

Let
\begin{equation}
 X^{(N)}=\sum_{j=1}^N2^{-j}U_j.
 \label{g2:eq:finite-dyadic-sum}
\end{equation}
The CDF of a sum of independent uniforms with widths $w_j$ is a truncated-power
spline.  For the dyadic widths, all subset sums form the uniform grid
$m/2^N$ and their inclusion--exclusion signs become the Thue--Morse signs.
Thus
\begin{equation}
 \Probability(X^{(N)}\le x)
 =\frac{2^{N(N+1)/2}}{N!}
 \sum_{m=0}^{2^N-1}(-1)^{\BinaryDigitSum(m)}
 \left(x-\frac{m}{2^N}\right)_+^N,
 \label{g2:eq:thue-morse-spline-cdf}
\end{equation}
where $\BinaryDigitSum(m)$ is the binary digit sum of $m$.  Differentiating gives the density with
power $N-1$ and factor $(N-1)!^{-1}$.

The Gaussian comb filter and the Thue--Morse spline are two different
projections of the same hierarchy of scales:

\begin{itemize}
\item the Lagrange filter keeps the scale index $j$ and uses only $n+1$
      multiplicative nodes $q^j$; Gaussian coefficients encode its exact
      moment cancellations;
\item the finite convolution expands all $2^N$ subsets of dyadic scales;
      subset sums become additive knots and parity produces Thue--Morse
      cancellation.
\end{itemize}

At $q=1/2$, the products $(1-2^{-j})$, the Mersenne factors $(2^j-1)$,
Gaussian coefficients at bases $1/2$ and $2$, and the binary subset lattice
all coexist.  This is why the Fabius/Rvachev system is not merely an example of
the general geometric theory: it is the point at which multiplicative
interpolation, additive splines, and binary arithmetic lock together exactly.


\section{Research frontiers}
\label{g2:sec:frontiers}

The problems below are ordered from sharply formulated local questions to
broader structural programs.  None is needed for the proved results.

\subsection{Asymptotics of the Fabius Newton coefficients}

\Cref{g2:conj:quarter-exponent} asks only for the leading quadratic exponent.  A
more informative target is a saddle-point expansion of
\eqref{g2:eq:fabius-Gaussian-transform} or
\eqref{g2:eq:newton-generating-factorization}.  The Gaussian coefficient favors
$j\approx n/2$, while $\mu_j/j!$ has its own lower-Lambert saddle and periodic
correction.  The interaction should produce a two-scale saddle problem.

\begin{problem}[Full coefficient asymptotic]
\label{g2:prob:full-A-asymptotic}
Determine functions $c_2,c_1,c_{\log}$ and a bounded oscillatory term
$\Psi$ such that
\begin{equation}
 \log|A_n|
 =c_2n^2+c_1n\log n+c_{\log}n
 +\Psi(\vartheta_n)+o(1),
 \label{g2:eq:full-A-asymptotic}
\end{equation}
with an explicit Lambert-type phase $\vartheta_n$.  Decide whether sign
changes and near-cancellations are controlled by the same periodic function
that occurs in the sharp Fabius endpoint asymptotic.
\end{problem}

A first rigorous step would be matching upper and lower bounds
$2^{n^2/4+o(n^2)}$.  The upper bound is already in
\Cref{g2:cor:coefficient-upper}; a lower bound must prevent exponential-scale
cancellation in the alternating central block.

\subsection{Boundary-crossing formulas for \texorpdfstring{$q>1/2$}{q greater than 1/2}}

\begin{problem}[The supercritical comb-moment law]
\label{g2:prob:qgt-half}
For $1/2<q<1$, derive a closed formula for $F_q(q^n)$ in terms of moments and
finite affine words generated by
\[
 x\mapsto x/q,
 \qquad
 x\mapsto(x-(1-q))/q.
\]
A satisfactory formula should specialize to \eqref{g2:eq:comb-moment} when no
word crosses the reflected boundary, and should make the breakpoints in $q$
explicit.
\end{problem}

The number of admissible boundary words may change when polynomial relations
among $q$-powers are crossed.  This suggests a piecewise-algebraic structure
related to beta-expansions, but with continuous uniform digits rather than
Bernoulli digits.

\subsection{Confluent limits and roots of unity}

\begin{problem}[$q\to0$ confluent calculus]
\label{g2:prob:q0-confluent}
Find a renormalized cardinal and Newton basis whose $q\to0$ limit separates the
infinitely many nodes collapsing at $0$.  Determine whether the correct limit
is a triangular mixture of value, derivative, and tail-quotient functionals.
\end{problem}

\begin{problem}[Cyclotomic geometric interpolation]
\label{g2:prob:roots-unity}
When $q$ is a primitive $m$th root of unity, replace repeated nodes by a
canonical Hermite data set and factor the resulting confluent Vandermonde
matrix into cyclotomic analogues of $P_q$ and $D_{a,q}$.  The denominator-free
operator
$\prod_{j=1}^n(E_q-q^j)$ should survive as a periodic difference operator even
when normalized Lagrange weights do not.
\end{problem}

These limits may connect the present theory to divided powers, quantum groups
at roots of unity, and cyclic $q$-difference equations.

\subsection{Regularization near Fourier zeros}

\begin{problem}[Filtered finite parts at sinc zeros]
\label{g2:prob:zero-regularization}
Let $\xi_0$ be a zero of $\FourierTwoPi{u_q}$ of known multiplicity $m$.  Construct
from $R_{n,q}$ and its first $m$ derivatives a sequence converging to the
finite part of $1/\FourierTwoPi{u_q}$ at $\xi_0$, with a bound uniform in a shrinking
neighborhood of the zero.
\end{problem}

Because every zero comes from explicit finite sinc factors, one can first
divide out the known local factor and then apply the geometric filter to the
nonvanishing tail.  The main issue is numerical cancellation when several
zero lattices overlap, as they do extensively at $q=1/2$.

\subsection{Exact Lebesgue functions below the finest node}

For $x=aq^\alpha$ with real $\alpha>n$, the cardinal signs are
$(-1)^{n-j}$, so the Lebesgue function is the explicit positive sum below.
The formula extends continuously to $\alpha=n$, where the target is the finest
node and the Lebesgue function equals $1$.
\begin{equation}
 \Lambda_n(aq^\alpha)=
 (q^{\alpha-n};q)_{n+1}
 \sum_{j=0}^n
 \frac{q^{\binom{n-j+1}{2}}}
 {(1-q^{\alpha-j})(q;q)_j(q;q)_{n-j}}.
 \label{g2:eq:below-comb-lebesgue}
\end{equation}
At $\alpha=\infty$ it reduces to \eqref{g1:eq:K-n}.  It is natural to ask
whether the finite sum has a single basic-hypergeometric closed form and
whether $\Lambda_n$ is monotone on $[0,aq^n]$.

\begin{conjecture}[Maximum below the comb]
\label{g2:conj:below-comb-lebesgue}
For $0<q<1$ and $a>0$,
\begin{equation}
 \max_{0\le x\le aq^n}\Lambda_n(x)=\Lambda_n(0)
 =\frac{(-q;q)_n}{(q;q)_n}.
 \label{g2:eq:below-comb-max}
\end{equation}
\end{conjecture}

Low-degree symbolic and numerical checks support the claim, but a proof
appears to require a variation-diminishing or total-positivity argument rather
than termwise monotonicity.

\subsection{Flatness-adapted approximation spaces}

Polynomial interpolation fails off the comb for a structural reason.  The
basis must be changed, not merely evaluated more accurately.  Classical
geometric-knot spline constructions~\cite{KocakPhillips1994} provide one
nearby approximation-theoretic model, while Rvachev atoms add exact
$C^\infty$ compact support and endpoint flatness.

\begin{problem}[Rvachev-preconditioned comb interpolation]
\label{g2:prob:rvachev-preconditioned}
Construct finite-dimensional spaces generated by scale-dependent Rvachev
atoms, for example
\[
 u(2^jx-r),\qquad 0\le j\le n,
\]
that interpolate the Fabius dyadic data and incorporate endpoint flatness.
Determine whether one can obtain convergence on $[0,1]$ while retaining exact
rational or algebraic coefficient formulas.
\end{problem}

The Gaussian--Appell decoder \eqref{gq:eq:cardinal-up-decoder} supplies an exact
preconditioner for each finite polynomial interpolant, but it cannot remove
the divergence because it reproduces the same polynomial.  A genuinely new
scheme must let the atom space grow in smoothness or scale, rather than merely
re-expressing the space of polynomials of degree at most $n$.

\begin{problem}[Growing Hermite data]
\label{g2:prob:growing-hermite}
At the tooth $2^{-j}$, include derivatives through an order $r_j$ that grows
with $j$.  Find the minimal growth of $r_j$ for which the Hermite interpolants
can converge to $F$ on a nontrivial interval.  Relate the threshold to the
sharp derivative and endpoint asymptotics of $F$.
\end{problem}

\subsection{Multivariate geometric meshes}

Schoenberg's one-dimensional problem has a triangular-mesh extension due to
Lee and Phillips~\cite{LeePhillips1988}.  The present operator and probability
viewpoints suggest several further versions.

\begin{problem}[Tensor and simplex filters]
\label{g2:prob:multivariate}
For nodes
$(a_1q_1^{j_1},\ldots,a_dq_d^{j_d})$, develop tensor and total-degree
Lagrange filters with explicit multivariate Mellin symbols.  On a simplex-like
geometric mesh, identify the correct multivariate Gaussian coefficients and
complete homogeneous residuals.  Apply them to joint geometric-uniform laws
and tensor products of up-functions.
\end{problem}

The tensor case is immediate algebraically but may suffer severe anisotropic
conditioning.  The simplex case should involve $q$-multinomial coefficients
and may admit a probabilistic interpretation through flags of coordinate
subspaces.

\subsection{Adaptive ratios and asymptotic optimization}

The exact error factors make ratio selection an optimization problem.  Small
$q$ gives powerful Gaussian suppression but spreads the first nodes; $q$ near
$1$ resolves a local germ but causes the condition number
\eqref{g1:eq:K-modular} to explode.

\begin{problem}[Optimal $q$ under noise]
\label{g2:prob:optimal-q}
Suppose the data $f(aq^k)$ have absolute or relative noise of known scale and
$f$ belongs to a specified analytic, Gevrey, or endpoint-asymptotic class.
Minimize a rigorous sum of truncation and amplified-noise bounds over both
$q$ and $n$.  Determine whether the optimum has a universal scaling law as the
noise tends to zero.
\end{problem}

For the Fabius function, the superflat signal and the large off-comb
coefficients make this a nonstandard regularization problem; the optimum for
evaluating $F(0)$ need not resemble the optimum for reconstructing $F(x)$ at
$x>0$.

\subsection{Formal verification roadmap}

The finite algebra is particularly well suited to Lean because most identities
can be stated over a commutative ring before any division is introduced.  A
productive dependency order is:

\begin{enumerate}
\item prove the denominator-free nodal product and geometric Vandermonde
      factorization;
\item define the Gaussian Pascal matrix and prove the inverse basis pair of
      \Cref{gq:thm:gaussian-pascal-basis};
\item derive the barycentric denominator and general-target cardinal after
      adding field hypotheses;
\item define the dilation polynomial and prove the full polynomial identity
      \eqref{g1:eq:row-polynomial};
\item obtain natural moment annihilation by evaluation, then extend to the
      analytic Mellin symbol in a complex-power layer;
\item formalize Jackson divided differences first for polynomials, then for
      functions with iterated $q$-derivatives;
\item prove the row-norm identity and Toeplitz regularity using summable
      reversed rows;
\item reuse the repository's geometric-uniform law to prove
      \eqref{g2:eq:comb-moment};
\item instantiate at $q=1/2$ and connect to existing Fabius moments and
      Mersenne products;
\item formalize the superflat divergence theorem using a local power-series
      contradiction;
\item feed the existing sinc-product theorem through the analytic filter to
      obtain \Cref{g2:thm:reciprocal-product}.
\end{enumerate}

The distinction between denominator-free polynomial statements and normalized
field statements should be maintained throughout.  It will make the later
root-of-unity and confluent developments substantially easier.
